%=============================================================================
% Chapter 4 Solutions
%=============================================================================
\newpage
\section*{Chapter 4 Solutions}
\addcontentsline{toc}{section}{Chapter 4 Solutions}

%-----------------------------------------------------------------------------
\subsection*{4.1 --- Control Signals for \texttt{and}}
\hypertarget{sol-4.1}{}
\markright{4.1}
\addcontentsline{toc}{subsection}{4.1 --- Control Signals for \texttt{and}}

Consider the instruction \texttt{and rd, rs1, rs2}, which computes
\texttt{Reg[rd] = Reg[rs1] AND Reg[rs2]}.  This is an R-type instruction.

\begin{center}
  \includegraphics[width=0.55\textwidth]{fig4_10_memory_rtype_datapath}
  \captionof{figure}{Figure 4.10 (reproduced): The datapath for the memory
    instructions and the R-type instructions.}
\end{center}

\phantomsection\addcontentsline{toc}{subsubsection}{4.1.1}
\medskip\noindent\textbf{4.1.1}
The control signals generated for \texttt{and} (an R-type instruction)
in Figure~4.10 are:

\begin{center}
\begin{tabular}{lcp{8cm}}
\toprule
\textbf{Signal} & \textbf{Value} & \textbf{Reason} \\
\midrule
RegWrite  & 1    & The result of the AND operation must be written back
                    to the register file (into \texttt{rd}). \\
ALUSrc    & 0    & The second ALU operand comes from a register
                    (\texttt{Read Data 2}), not from the sign-extended
                    immediate. \\
ALU operation & 0000 & The ALU must perform the bitwise AND operation. \\
MemRead   & 0    & No data is read from memory. \\
MemWrite  & 0    & No data is written to memory. \\
MemtoReg  & 0    & The value written back to the register file comes
                    from the ALU result, not from data memory. \\
\bottomrule
\end{tabular}
\end{center}

\phantomsection\addcontentsline{toc}{subsubsection}{4.1.2}
\medskip\noindent\textbf{4.1.2}
The resources that perform a \emph{useful} function for \texttt{and}:

\begin{itemize}
  \item \textbf{Register file} --- reads \texttt{rs1} and \texttt{rs2}
        (Read Data 1 and Read Data 2), and writes the result into
        \texttt{rd} (RegWrite = 1).
  \item \textbf{ALU} --- computes the bitwise AND of its two inputs and
        produces the result.
  \item \textbf{ALUSrc mux} --- selects the register value
        (Read Data 2) as the second ALU input (ALUSrc = 0).
  \item \textbf{MemtoReg mux} --- selects the ALU result as the data
        to be written back to the register file (MemtoReg = 0).
\end{itemize}

\phantomsection\addcontentsline{toc}{subsubsection}{4.1.3}
\medskip\noindent\textbf{4.1.3}

\textbf{Resources that produce no output:}

\begin{itemize}
  \item \textbf{Data memory} --- MemRead = 0 and MemWrite = 0, so no
        read or write operation occurs.  The data memory is completely
        idle for this instruction.
\end{itemize}

\textbf{Resources that produce output that is not used:}

\begin{itemize}
  \item \textbf{Sign-extend unit} --- The sign-extend unit is
        combinational and always produces a 32-bit sign-extended version
        of bits [15:0] of the instruction.  However, because ALUSrc = 0,
        the ALUSrc mux selects Read Data 2 instead, so the sign-extended
        output is ignored.
  \item \textbf{ALU Zero output} --- The ALU always produces the Zero
        flag, but for \texttt{and} this signal is not used (it is only
        relevant for branch instructions).
\end{itemize}

%-----------------------------------------------------------------------------
\subsection*{4.2 --- Don't Cares in Control}
\hypertarget{sol-4.2}{}
\markright{4.2}
\addcontentsline{toc}{subsection}{4.2 --- Don't Cares in Control}

\begin{center}
  \includegraphics[width=0.65\textwidth]{fig4_18_control_signal_truth_table}
  \captionof{figure}{Figure 4.18 (reproduced): The control function for the
    simple single-cycle implementation.}
\end{center}

Figure~4.18 contains four ``X'' (don't care) entries.  They all arise
from the same underlying reason: \textbf{these signals only matter when
the register file is written to (RegWrite = 1)}.  For instructions with
RegWrite = 0, no value is written to any register, so it does not matter
which register is selected as the destination or what data is routed to
the write port.

\medskip
\begin{center}
\begin{tabular}{llp{8.5cm}}
\toprule
\textbf{Instruction} & \textbf{Don't care} & \textbf{Explanation} \\
\midrule
\texttt{sw} & RegDst = X &
  \texttt{sw} does not write to the register file (RegWrite = 0).
  RegDst controls whether the write-register number comes from
  \texttt{rt} (bits [20:16]) or \texttt{rd} (bits [15:11]).
  Since no register is written, the mux selection is irrelevant. \\[6pt]
\texttt{sw} & MemtoReg = X &
  MemtoReg selects whether the ALU result or the data-memory output is
  sent to the register file's write-data port.  Since \texttt{sw} has
  RegWrite = 0, nothing is written regardless of what appears on the
  write-data bus. \\[6pt]
\texttt{beq} & RegDst = X &
  \texttt{beq} does not write to the register file (RegWrite = 0).
  The write-register selection is therefore irrelevant, same as
  for \texttt{sw}. \\[6pt]
\texttt{beq} & MemtoReg = X &
  As with \texttt{sw}, since RegWrite = 0 for \texttt{beq}, it does not
  matter which value (ALU result or memory read data) is presented to the
  register file's write-data input. \\
\bottomrule
\end{tabular}
\end{center}

\medskip\noindent
\textbf{Summary:} In both \texttt{sw} and \texttt{beq}, RegWrite = 0,
which disables the register file write.  This makes \emph{both} RegDst
and MemtoReg don't cares, because:
\begin{itemize}
  \item RegDst feeds the write-register mux, which selects \emph{where}
        to write --- irrelevant if nothing is written.
  \item MemtoReg feeds the write-data mux, which selects \emph{what}
        to write --- also irrelevant if nothing is written.
\end{itemize}
R-format and \texttt{lw} both have RegWrite = 1, so they \emph{must}
set RegDst and MemtoReg to specific values to ensure the correct
register receives the correct data.

%-----------------------------------------------------------------------------
\subsection*{4.3 --- Instruction Mix Resource Usage}
\hypertarget{sol-4.3}{}
\markright{4.3}
\addcontentsline{toc}{subsection}{4.3 --- Instruction Mix Resource Usage}

Given instruction mix: R-type 24\%, I-type (non-lw) 28\%, Load 25\%,
Store 10\%, Branch 11\%, Jump 2\%.

\phantomsection\addcontentsline{toc}{subsubsection}{4.3.1}
\medskip\noindent\textbf{4.3.1}
The data memory is used by instructions that read from or write to memory:
\[
  \text{Fraction using data memory} = \text{Load} + \text{Store}
  = 25\% + 10\% = \boxed{35\%}
\]

\phantomsection\addcontentsline{toc}{subsubsection}{4.3.2}
\medskip\noindent\textbf{4.3.2}
The sign-extend unit is used by all instructions that have an immediate
field: I-type (non-lw), load, store, and branch.
\[
  \text{Fraction using sign extend} = 28\% + 25\% + 10\% + 11\%
  = \boxed{74\%}
\]

\noindent Note: R-type (24\%) and jump (2\%) do not use the sign-extended
immediate, totaling 26\%.  As a check: $74\% + 26\% = 100\%$. \checkmark

\phantomsection\addcontentsline{toc}{subsubsection}{4.3.3}
\medskip\noindent\textbf{4.3.3}
The sign-extend unit is a purely combinational block; it \emph{always}
produces an output every cycle, regardless of the instruction type.  During
cycles when its output is not needed (R-type and jump instructions), the
sign-extended value is still computed, but it is simply \emph{not selected}
by the ALUSrc mux.

\begin{itemize}
  \item For R-type instructions, ALUSrc = 0, so the mux selects the
        register value (Read Data 2) rather than the sign-extended immediate.
  \item For jump instructions, the ALU result is ignored entirely.
\end{itemize}

The sign-extend unit consumes a small amount of power during these cycles
but does not affect correctness---the unwanted output is gated off by the
mux.

\begin{center}
  \includegraphics[width=0.45\textwidth]{fig4_24_datapath_with_jump}
  \captionof{figure}{Figure 4.24: Jump uses a dedicated path
    (instruction [25:0] shifted left 2, concatenated with PC+4[31:28])
    to determine the next PC, bypassing the ALU entirely.}
\end{center}

%-----------------------------------------------------------------------------
\subsection*{4.4 --- Stuck-at-0 Fault Analysis}
\hypertarget{sol-4.4}{}
\markright{4.4}
\addcontentsline{toc}{subsection}{4.4 --- Stuck-at-0 Fault Analysis}

\phantomsection\addcontentsline{toc}{subsubsection}{4.4.1}
\medskip\noindent\textbf{4.4.1}
\textbf{RegDst stuck at 0:}

In the normal control table, RegDst = 1 only for R-type instructions.
When RegDst = 1, the write-register number comes from bits [15:11] (the
\texttt{rd} field).  When RegDst = 0, it comes from bits [20:16] (the
\texttt{rt} field).

If RegDst is stuck at 0:
\begin{itemize}
  \item \textbf{R-type instructions fail.}  They need RegDst = 1 to write
        to the \texttt{rd} register.  With RegDst stuck at 0, the result is
        written to \texttt{rt} instead of \texttt{rd}.  For example,
        \texttt{add \$s0, \$s1, \$s2} would write the sum to \texttt{\$s2}
        (rt) instead of \texttt{\$s0} (rd).
  \item \textbf{lw instructions are unaffected.}  They already require
        RegDst = 0 (write to \texttt{rt}).
  \item \textbf{sw, beq, j are unaffected.}  RegDst is a don't-care for
        these instructions because they do not write to the register file
        (RegWrite = 0).
\end{itemize}

\phantomsection\addcontentsline{toc}{subsubsection}{4.4.2}
\medskip\noindent\textbf{4.4.2}
\textbf{ALUSrc stuck at 0:}

ALUSrc controls the second input to the ALU.  When ALUSrc = 0, the ALU
receives Read Data 2 (a register value).  When ALUSrc = 1, it receives the
sign-extended immediate.

If ALUSrc is stuck at 0:
\begin{itemize}
  \item \textbf{lw and sw instructions fail.}  Both need ALUSrc = 1 so the
        ALU computes \texttt{base + offset} using the sign-extended
        immediate.  With ALUSrc = 0, the ALU receives the value of
        \texttt{rt} instead of the immediate, producing an incorrect memory
        address.
  \item \textbf{I-type ALU instructions (e.g., addi, andi, ori) also fail.}
        They need ALUSrc = 1 to use the immediate operand.  With ALUSrc = 0,
        the ALU uses a register value instead.
  \item \textbf{R-type instructions are unaffected.}  They already require
        ALUSrc = 0.
  \item \textbf{beq is unaffected.}  It requires ALUSrc = 0 (subtracts
        two registers to compare).
  \item \textbf{j is unaffected.}  ALUSrc is a don't-care for jumps.
\end{itemize}

%-----------------------------------------------------------------------------
\subsection*{4.5 --- Trace Instruction Through Datapath}
\hypertarget{sol-4.5}{}
\markright{4.5}
\addcontentsline{toc}{subsection}{4.5 --- Trace Instruction Through Datapath}

Instruction: \texttt{add \$s0, \$s1, \$s2}

Encoding (R-type):
\begin{verbatim}
    opcode = 000000  (R-type)
    rs     = 10001   ($s1 = register 17)
    rt     = 10010   ($s2 = register 18)
    rd     = 10000   ($s0 = register 16)
    shamt  = 00000
    funct  = 100000  (add)
\end{verbatim}

\phantomsection\addcontentsline{toc}{subsubsection}{4.5.1}
\medskip\noindent\textbf{4.5.1}
The control unit decodes opcode = \texttt{000000} (R-type) and generates:

\begin{center}
\begin{tabular}{lc}
\toprule
\textbf{Signal} & \textbf{Value} \\
\midrule
RegDst   & 1 \\
ALUSrc   & 0 \\
MemtoReg & 0 \\
RegWrite & 1 \\
MemRead  & 0 \\
MemWrite & 0 \\
Branch   & 0 \\
ALUOp    & 10 \\
Jump     & 0 \\
\bottomrule
\end{tabular}
\end{center}

The ALU control unit receives ALUOp = 10 and funct = \texttt{100000},
producing ALU control = \texttt{0010} (addition).

\phantomsection\addcontentsline{toc}{subsubsection}{4.5.2}
\medskip\noindent\textbf{4.5.2}
The new PC value is:
\[
  \text{PC}_{\text{new}} = \text{PC} + 4
\]

Path: The current PC feeds the instruction memory \emph{and} the top adder,
which computes PC + 4.  Since Branch = 0 and Jump = 0, the branch-target
mux selects PC + 4, and the jump mux also selects that value.  The result
is written to the PC register on the next rising clock edge.

No branch is taken because Branch AND Zero = $0 \text{ AND } (\cdot) = 0$.

\phantomsection\addcontentsline{toc}{subsubsection}{4.5.3}
\medskip\noindent\textbf{4.5.3}
There are three main muxes in the datapath:

\begin{center}
\begin{tabular}{lcccc}
\toprule
\textbf{Mux} & \textbf{Input 0} & \textbf{Input 1} & \textbf{Select} & \textbf{Output} \\
\midrule
RegDst mux   & rt = 18 & rd = 16 & 1 & 16 (\texttt{\$s0}) \\
ALUSrc mux   & Read Data 2 = [\texttt{\$s2}] & Sign-ext imm & 0 & [\texttt{\$s2}] \\
MemtoReg mux & ALU result & Mem read data & 0 & ALU result \\
\bottomrule
\end{tabular}
\end{center}

Additionally, the PC-source mux:
\begin{itemize}
  \item Input 0 = PC + 4, Input 1 = branch target (PC + 4 + offset$\times$4).
  \item Select = Branch AND Zero = 0.  Output = PC + 4.
\end{itemize}

And the Jump mux:
\begin{itemize}
  \item Input 0 = output of PC-source mux = PC + 4.
        Input 1 = jump target.
  \item Select = Jump = 0.  Output = PC + 4.
\end{itemize}

\phantomsection\addcontentsline{toc}{subsubsection}{4.5.4}
\medskip\noindent\textbf{4.5.4}
\textbf{Register File read outputs:}
\begin{itemize}
  \item Read Register 1 = rs = 17 $\Rightarrow$ Read Data 1 = contents of \texttt{\$s1}
  \item Read Register 2 = rt = 18 $\Rightarrow$ Read Data 2 = contents of \texttt{\$s2}
\end{itemize}

\textbf{ALU inputs:}
\begin{itemize}
  \item Input A = Read Data 1 = contents of \texttt{\$s1}
  \item Input B = output of ALUSrc mux = Read Data 2 = contents of \texttt{\$s2}
        (since ALUSrc = 0)
  \item ALU control = \texttt{0010} (add)
  \item ALU result = \texttt{\$s1} + \texttt{\$s2}
  \item Zero flag = 1 if the result is zero, 0 otherwise (not used by this
        instruction)
\end{itemize}

\phantomsection\addcontentsline{toc}{subsubsection}{4.5.5}
\medskip\noindent\textbf{4.5.5}
The Register File write inputs are:

\begin{center}
\begin{tabular}{lc}
\toprule
\textbf{Input} & \textbf{Value} \\
\midrule
Write Register & 16 (from RegDst mux, selecting rd = \texttt{\$s0}) \\
Write Data     & ALU result = \texttt{\$s1} + \texttt{\$s2} \\
               & (from MemtoReg mux, selecting ALU result) \\
RegWrite       & 1 (write is enabled) \\
\bottomrule
\end{tabular}
\end{center}

At the next rising clock edge, the value \texttt{\$s1} + \texttt{\$s2} is
written into register \texttt{\$s0} (register 16).
%-----------------------------------------------------------------------------
\subsection*{4.6 --- Adding I-type Instructions}
\hypertarget{sol-4.6}{}
\markright{4.6}
\addcontentsline{toc}{subsection}{4.6 --- Adding I-type Instructions}

Section 4.1 discusses R-type, load, store, branch, and jump instructions
but omits I-type ALU instructions such as \texttt{addi}, \texttt{andi},
\texttt{ori}, and \texttt{slti}.

\phantomsection\addcontentsline{toc}{subsubsection}{4.6.1}
\medskip\noindent\textbf{4.6.1}
No new logic blocks are needed for most I-type ALU instructions.  The
existing single-cycle datapath (Figure~4.17) already contains every
component required:

\begin{itemize}
  \item The \textbf{sign-extend unit} produces the 32-bit immediate from the
        16-bit field in the instruction.
  \item The \textbf{ALUSrc mux} can already select the sign-extended
        immediate as the second ALU input.
  \item The \textbf{RegDst mux} can already select \texttt{rt}
        (bits [20:16]) as the write register.
  \item The ALU already supports add, subtract, AND, OR, and set-on-less-than
        operations.
\end{itemize}

\noindent\textbf{Exception:} For \texttt{andi} and \texttt{ori}, the
immediate must be \emph{zero-extended} rather than sign-extended.  To support
this we need either:
\begin{enumerate}
  \item A \textbf{zero-extend unit} and a \textbf{mux} to select between
        sign-extension and zero-extension, controlled by a new signal
        (e.g.\ \texttt{ExtOp}), or
  \item Modify the existing sign-extend unit to accept a control input that
        chooses between sign- and zero-extension.
\end{enumerate}

\phantomsection\addcontentsline{toc}{subsubsection}{4.6.2}
\medskip\noindent\textbf{4.6.2}
Control signal values for representative I-type instructions:

\begin{center}
\begin{tabular}{lccccccccc}
\toprule
\textbf{Signal} & \textbf{addi} & \textbf{andi} & \textbf{ori} & \textbf{slti} \\
\midrule
RegDst   & 0  & 0  & 0  & 0  \\
ALUSrc   & 1  & 1  & 1  & 1  \\
MemtoReg & 0  & 0  & 0  & 0  \\
RegWrite & 1  & 1  & 1  & 1  \\
MemRead  & 0  & 0  & 0  & 0  \\
MemWrite & 0  & 0  & 0  & 0  \\
Branch   & 0  & 0  & 0  & 0  \\
ALUOp    & 00 & 11 & 11 & 11 \\
Jump     & 0  & 0  & 0  & 0  \\
\bottomrule
\end{tabular}
\end{center}

Reasoning for each signal:
\begin{itemize}
  \item \textbf{RegDst = 0}: The destination register is \texttt{rt}
        (bits [20:16]), not \texttt{rd} (bits [15:11]).
  \item \textbf{ALUSrc = 1}: The second ALU operand is the immediate field
        (sign- or zero-extended), not a register value.
  \item \textbf{MemtoReg = 0}: The result written to the register file comes
        from the ALU output, not from data memory.
  \item \textbf{RegWrite = 1}: The instruction writes a result to a register.
  \item \textbf{MemRead = 0}: No data memory read is performed.
  \item \textbf{MemWrite = 0}: No data memory write is performed.
  \item \textbf{Branch = 0}: This is not a branch instruction.
  \item \textbf{ALUOp}: For \texttt{addi}, ALUOp = 00 (addition).  For
        \texttt{andi}, \texttt{ori}, and \texttt{slti}, ALUOp must encode a
        value that causes the ALU control to select the correct operation
        (AND, OR, or SLT).  Since the existing 2-bit ALUOp only encodes
        R-type (10), load/store (00), and branch (01), a new encoding such
        as 11 can be used, with the \texttt{funct}-equivalent operation
        determined by the opcode bits.
  \item \textbf{Jump = 0}: This is not a jump instruction.
\end{itemize}

\noindent There are no true ``don't care'' signals for I-type ALU
instructions---every control signal has a well-defined required value.

%-----------------------------------------------------------------------------
\subsection*{4.7 --- Instruction Latencies}
\hypertarget{sol-4.7}{}
\markright{4.7}
\addcontentsline{toc}{subsection}{4.7 --- Instruction Latencies}

We trace the critical path through the single-cycle datapath for each
instruction type using the given latencies:

\begin{center}
\begin{tabular}{ccccccccc}
\toprule
\textbf{I-Mem} & \textbf{Mux} & \textbf{ALU} & \textbf{Adder} &
\textbf{Gate} & \textbf{Reg} & \textbf{Reg File} &
\textbf{Sign ext.} & \textbf{Ctrl} \\
\midrule
200ps & 30ps & 120ps & 150ps & 5ps & 90ps & 150ps & 10ps & 50ps \\
\bottomrule
\end{tabular}
\end{center}

\noindent Note: ``Reg'' (90\,ps) is the register setup time (data must be
stable before the clock edge).  ``Reg File'' (150\,ps) is the register read
time (output delay after clock edge).  We assume data memory has the same
latency as instruction memory (200\,ps).

\phantomsection\addcontentsline{toc}{subsubsection}{4.7.1}
\medskip\noindent\textbf{4.7.1}
\textbf{R-type critical path:}
\[
  \underbrace{200}_{\text{I-Mem}} +
  \underbrace{\max(150,\,50)}_{\text{max(Reg read, Ctrl)}} +
  \underbrace{30}_{\text{Mux (ALUSrc)}} +
  \underbrace{120}_{\text{ALU}} +
  \underbrace{30}_{\text{Mux (MemtoReg)}} +
  \underbrace{90}_{\text{Reg setup}}
\]
\[
  = 200 + 150 + 30 + 120 + 30 + 90 = \boxed{620\text{\,ps}}
\]

The parallel paths after instruction fetch are:
\begin{itemize}
  \item Control path: Ctrl (50\,ps) $\to$ mux select signals
  \item Data path: Reg File read (150\,ps)
\end{itemize}
Since $150 > 50$, the register file read is on the critical path.

\phantomsection\addcontentsline{toc}{subsubsection}{4.7.2}
\medskip\noindent\textbf{4.7.2}
\textbf{\texttt{lw} critical path:}
\[
  \underbrace{200}_{\text{I-Mem}} +
  \underbrace{150}_{\text{Reg read}} +
  \underbrace{30}_{\text{Mux (ALUSrc)}} +
  \underbrace{120}_{\text{ALU}} +
  \underbrace{200}_{\text{D-Mem}} +
  \underbrace{30}_{\text{Mux (MemtoReg)}} +
  \underbrace{90}_{\text{Reg setup}}
\]
\[
  = 200 + 150 + 30 + 120 + 200 + 30 + 90 = \boxed{820\text{\,ps}}
\]

This is longer than R-type because \texttt{lw} must also access data memory
(200\,ps) between the ALU and the write-back mux.

\phantomsection\addcontentsline{toc}{subsubsection}{4.7.3}
\medskip\noindent\textbf{4.7.3}
\textbf{I-type (non-load) critical path} (e.g.\ \texttt{addi}):

The path is identical to R-type.  Although the immediate passes through the
sign-extend unit (10\,ps) and then the ALUSrc mux, this path
(Ctrl 50\,ps + Sign-ext 10\,ps = 60\,ps) is still shorter than the register
read path (150\,ps).  Therefore:
\[
  200 + 150 + 30 + 120 + 30 + 90 = \boxed{620\text{\,ps}}
\]

\phantomsection\addcontentsline{toc}{subsubsection}{4.7.4}
\medskip\noindent\textbf{4.7.4}
The minimum clock period must accommodate the slowest instruction.  We also
check the remaining instruction types:

\textbf{\texttt{sw} critical path:}
\[
  200\;(\text{I-Mem}) + 150\;(\text{Reg read}) + 30\;(\text{Mux}) +
  120\;(\text{ALU}) + 200\;(\text{D-Mem write}) = 700\text{\,ps}
\]
(\texttt{sw} does not write back to the register file.)

\textbf{\texttt{beq} critical path:}
\[
  200\;(\text{I-Mem}) + 150\;(\text{Reg read}) + 30\;(\text{Mux}) +
  120\;(\text{ALU}) + 5\;(\text{AND gate}) + 30\;(\text{Mux for PC})
  = 535\text{\,ps}
\]

\textbf{\texttt{j} critical path:}
The jump target is computed from the instruction word (shift and concatenate),
requiring no ALU or memory access:
$200 + 30 = 230$\,ps.

\medskip
\begin{center}
\begin{tabular}{lc}
\toprule
\textbf{Instruction Type} & \textbf{Latency} \\
\midrule
R-type      & 620\,ps \\
I-type (non-load) & 620\,ps \\
\texttt{lw} & 820\,ps \\
\texttt{sw} & 700\,ps \\
\texttt{beq} & 535\,ps \\
\texttt{j}  & 230\,ps \\
\bottomrule
\end{tabular}
\end{center}

The minimum clock period is determined by the longest latency:
\[
  \text{Minimum clock period} = \max(620, 620, 820, 700, 535, 230)
  = \boxed{820\text{\,ps}}
\]

%-----------------------------------------------------------------------------
\subsection*{4.8 --- Variable Clock Cycle Speedup}
\hypertarget{sol-4.8}{}
\markright{4.8}
\addcontentsline{toc}{subsection}{4.8 --- Variable Clock Cycle Speedup}

If each instruction type could use its own (variable) clock cycle time, the
average clock cycle would be the weighted average over the instruction mix:

\begin{center}
\begin{tabular}{lcc}
\toprule
\textbf{Instruction Type} & \textbf{Latency} & \textbf{Fraction} \\
\midrule
R-type            & 620\,ps & 24\% \\
I-type (non-lw)   & 620\,ps & 28\% \\
Load (\texttt{lw})& 820\,ps & 25\% \\
Store (\texttt{sw})& 700\,ps & 10\% \\
Branch            & 535\,ps & 11\% \\
Jump              & 230\,ps & 2\% \\
\bottomrule
\end{tabular}
\end{center}

\[
  \text{Avg cycle} = 0.24(620) + 0.28(620) + 0.25(820) + 0.10(700) +
                     0.11(535) + 0.02(230)
\]
\begin{align*}
  &= 148.8 + 173.6 + 205.0 + 70.0 + 58.85 + 4.6 \\
  &= 660.85\text{\,ps}
\end{align*}

The fixed-clock CPU must use $820$\,ps (the \texttt{lw} latency) for every
instruction.  For a program of $N$ instructions:
\[
  \text{Speedup} = \frac{820}{660.85} \approx \boxed{1.24}
\]

The variable-clock CPU is approximately 24\% faster, primarily because it
avoids paying the full \texttt{lw} latency for the 75\% of instructions that
do not access data memory.

%-----------------------------------------------------------------------------
\subsection*{4.9 --- Multiplier Addition to CPU}
\hypertarget{sol-4.9}{}
\markright{4.9}
\addcontentsline{toc}{subsection}{4.9 --- Multiplier Addition to CPU}

Adding a multiplier increases the ALU latency by 300\,ps (from 120\,ps to
420\,ps) but reduces the instruction count by 5\%.

\phantomsection\addcontentsline{toc}{subsubsection}{4.9.1}
\medskip\noindent\textbf{4.9.1}
The clock cycle time is determined by the critical path (\texttt{lw}):

\textbf{Without improvement:}
\[
  200 + 150 + 30 + 120 + 200 + 30 + 90 = \boxed{820\text{\,ps}}
\]

\textbf{With improvement} (ALU latency = $120 + 300 = 420$\,ps):
\[
  200 + 150 + 30 + 420 + 200 + 30 + 90 = \boxed{1120\text{\,ps}}
\]

\phantomsection\addcontentsline{toc}{subsubsection}{4.9.2}
\medskip\noindent\textbf{4.9.2}
Let $N$ be the original instruction count.  The new instruction count is
$0.95N$.

\[
  \text{Speedup} = \frac{\text{Old execution time}}{\text{New execution time}}
  = \frac{820 \times N}{1120 \times 0.95N}
  = \frac{820}{1064} \approx \boxed{0.771}
\]

Since the speedup is less than 1, this is actually a \textbf{slowdown} of
approximately $1/0.771 \approx 1.30\times$.  The 36.6\% increase in clock
cycle time far outweighs the 5\% reduction in instruction count.  This
improvement should \textbf{not} be made.

\phantomsection\addcontentsline{toc}{subsubsection}{4.9.3}
\medskip\noindent\textbf{4.9.3}
Let $x$ be the additional latency added to the ALU (so the new ALU latency is
$120 + x$).  For the improvement to yield a net performance gain:
\[
  \text{Old time} > \text{New time}
\]
\[
  820 \times N > \bigl(200 + 150 + 30 + (120 + x) + 200 + 30 + 90\bigr)
  \times 0.95N
\]
\[
  820 > (820 + x) \times 0.95
\]
\[
  \frac{820}{0.95} > 820 + x
\]
\[
  863.16 > 820 + x
\]
\[
  x < 43.16\text{\,ps}
\]

The new ALU can be at most $120 + 43 = 163$\,ps, which means the multiplier
can add at most approximately \boxed{43\text{\,ps}} to the ALU latency and
still yield improved performance.

%-----------------------------------------------------------------------------
\subsection*{4.10 --- Register File Expansion Cost}
\hypertarget{sol-4.10}{}
\markright{4.10}
\addcontentsline{toc}{subsection}{4.10 --- Register File Expansion Cost}

Doubling the register file from 32 to 64 registers has three effects:
\begin{enumerate}
  \item Reduces the number of \texttt{lw} and \texttt{sw} instructions
        executed by 12\%.
  \item Increases register file read latency from 150\,ps to 160\,ps.
  \item Doubles the register file cost from 200 to 400.
\end{enumerate}

\phantomsection\addcontentsline{toc}{subsubsection}{4.10.1}
\medskip\noindent\textbf{4.10.1}
\textbf{Old clock cycle} (determined by \texttt{lw} critical path, Reg File
read = 150\,ps):
\[
  200 + 150 + 30 + 120 + 200 + 30 + 90 = 820\text{\,ps}
\]

\textbf{New clock cycle} (Reg File read = 160\,ps):
\[
  200 + 160 + 30 + 120 + 200 + 30 + 90 = 830\text{\,ps}
\]

From the instruction mix, \texttt{lw} and \texttt{sw} together account for
$25\% + 10\% = 35\%$ of instructions.  Reducing these by 12\% gives a new
instruction count:
\[
  N_{\text{new}} = N \times \bigl(1 - 0.12 \times 0.35\bigr)
  = N \times (1 - 0.042) = 0.958N
\]

Execution times:
\begin{align*}
  T_{\text{old}} &= 820 \times N \\
  T_{\text{new}} &= 830 \times 0.958N = 795.14N
\end{align*}

\[
  \text{Speedup} = \frac{820N}{795.14N}
  = \frac{820}{795.14} \approx \boxed{1.031}
\]

The improvement yields approximately a \textbf{3.1\% speedup}.

\phantomsection\addcontentsline{toc}{subsubsection}{4.10.2}
\medskip\noindent\textbf{4.10.2}
\textbf{Performance change:} $+3.1\%$.

\textbf{Cost change:} The register file cost doubles from 200 to 400,
an increase of $\Delta = 200$.  If the total CPU cost is approximately
1000 (as given in the textbook's cost table), then:
\[
  \text{Cost increase} = \frac{400 - 200}{1000} = 20\%
\]

The cost/performance ratio:
\[
  \frac{\text{Cost increase}}{\text{Performance increase}}
  = \frac{20\%}{3.1\%} \approx 6.5
\]

The cost increases by about 6.5$\times$ more (in percentage terms) than the
performance improves.  This is a \textbf{poor} cost/performance tradeoff.

\phantomsection\addcontentsline{toc}{subsubsection}{4.10.3}
\medskip\noindent\textbf{4.10.3}
\textbf{When more registers make sense:}
\begin{itemize}
  \item In \textbf{embedded systems} or systems with slow memory where
        reducing memory accesses is critical.  If the program is
        register-starved (many spills to stack), doubling registers could
        eliminate far more than 12\% of loads/stores, making the tradeoff
        worthwhile.
  \item In \textbf{deeply pipelined or out-of-order processors} where
        register renaming benefits from a larger physical register file,
        and where the register file latency increase is hidden by other
        pipeline stages.
\end{itemize}

\textbf{When more registers do not make sense:}
\begin{itemize}
  \item In \textbf{cost-sensitive designs} (e.g., low-power
        microcontrollers) where the 20\% cost increase is not justified by
        the modest 3\% performance gain.
  \item In \textbf{workloads with good cache behavior} and few register
        spills, where the 12\% reduction in \texttt{lw}/\texttt{sw} is
        optimistic and the actual benefit would be even smaller.
  \item When the increased register file size forces \textbf{wider
        instruction encodings} (more bits for register specifiers), which
        increases code size and may degrade instruction cache performance.
\end{itemize}
%-----------------------------------------------------------------------------
\subsection*{4.11 --- Load With Increment Instruction}
\hypertarget{sol-4.11}{}
\markright{4.11}
\addcontentsline{toc}{subsection}{4.11 --- Load With Increment Instruction}

The proposed instruction is:
\begin{verbatim}
    lwi $t1, 0($s1)    # $t1 = Mem[$s1]; $s1 = $s1 + 4
\end{verbatim}

This instruction performs two operations simultaneously: (1) load a word from
memory at address \texttt{\$s1} into \texttt{\$t1}, and (2) increment the base
register \texttt{\$s1} by 4.

\phantomsection\addcontentsline{toc}{subsubsection}{4.11.1}
\medskip\noindent\textbf{4.11.1}
We need \textbf{one new adder} to compute \texttt{\$s1 + 4}. The existing ALU
is already occupied computing the memory address \texttt{\$s1 + offset}
(which in this case is \texttt{\$s1 + 0}). Since both computations must happen
in the same cycle, we cannot reuse the ALU for the increment. The new adder
takes the value of \texttt{\$s1} (from Read Data~1 of the Register File) and
the constant~4 as inputs, and produces \texttt{\$s1 + 4} as output.

\phantomsection\addcontentsline{toc}{subsubsection}{4.11.2}
\medskip\noindent\textbf{4.11.2}
The \textbf{Register File} must be modified to support \textbf{two simultaneous
writes} in a single clock cycle. In the original design the Register File has
one write port (Write Register, Write Data, RegWrite). For \texttt{lwi}, we
must write both:
\begin{itemize}
  \item \texttt{\$t1} $\leftarrow$ the loaded data from memory, and
  \item \texttt{\$s1} $\leftarrow$ \texttt{\$s1 + 4} from the new adder.
\end{itemize}
This requires adding a \textbf{second write port} (Write Register~2, Write
Data~2, RegWrite2) to the Register File.

\phantomsection\addcontentsline{toc}{subsubsection}{4.11.3}
\medskip\noindent\textbf{4.11.3}
The following new data paths are needed:
\begin{enumerate}
  \item A path carrying the \textbf{constant~4} to one input of the new adder.
  \item A path from \textbf{Read Data~1} (the value of \texttt{\$s1}) to the
        other input of the new adder.
  \item A path from the \textbf{new adder's output} to the second write data
        port (Write Data~2) of the Register File.
  \item A path from the \textbf{rs field} of the instruction (bits [25:21]) to
        the second write register port (Write Register~2) of the Register File,
        so the incremented value is written back to \texttt{\$s1}.
\end{enumerate}

\phantomsection\addcontentsline{toc}{subsubsection}{4.11.4}
\medskip\noindent\textbf{4.11.4}
The following new control signals are needed:
\begin{itemize}
  \item \textbf{RegWrite2}: Enables the second write port of the Register File.
        Asserted only for \texttt{lwi}; deasserted for all other instructions.
  \item \textbf{WriteRegSrc2}: Selects which register number is sent to Write
        Register~2 (the \texttt{rs} field for \texttt{lwi}).
\end{itemize}
No additional ALU control is needed because the new adder is a dedicated
add-by-4 unit (hardwired), not a general-purpose ALU.

%-----------------------------------------------------------------------------
\subsection*{4.12 --- Swap Instruction}
\hypertarget{sol-4.12}{}
\markright{4.12}
\addcontentsline{toc}{subsection}{4.12 --- Swap Instruction}

The proposed instruction is:
\begin{verbatim}
    swap $s1, $s2       # temp = $s1; $s1 = $s2; $s2 = temp
\end{verbatim}

This instruction atomically exchanges the contents of two registers.

\phantomsection\addcontentsline{toc}{subsubsection}{4.12.1}
\medskip\noindent\textbf{4.12.1}
\textbf{No new functional blocks} are needed. The existing Register File can
already read both registers simultaneously (via Read Register~1 and Read
Register~2). No ALU computation or memory access is required --- the swap is
purely a register-to-register data movement.

\phantomsection\addcontentsline{toc}{subsubsection}{4.12.2}
\medskip\noindent\textbf{4.12.2}
The \textbf{Register File} must be modified to support \textbf{two simultaneous
writes}. In the original single-cycle datapath, the Register File has only one
write port. For \texttt{swap}, we must write:
\begin{itemize}
  \item \texttt{\$s1} $\leftarrow$ old value of \texttt{\$s2} (from Read
        Data~2), and
  \item \texttt{\$s2} $\leftarrow$ old value of \texttt{\$s1} (from Read
        Data~1).
\end{itemize}
This requires a \textbf{second write port} (Write Register~2, Write Data~2,
RegWrite2).

\phantomsection\addcontentsline{toc}{subsubsection}{4.12.3}
\medskip\noindent\textbf{4.12.3}
The following new data paths are needed:
\begin{enumerate}
  \item \textbf{Read Data~2 $\to$ Write Data~1}: The value read from
        \texttt{\$s2} must be routed to the first write data port so it can
        be written into \texttt{\$s1}. (A mux is needed at Write Data~1 to
        select between the normal write-back sources and Read Data~2.)
  \item \textbf{Read Data~1 $\to$ Write Data~2}: The value read from
        \texttt{\$s1} must be routed to the second write data port so it can
        be written into \texttt{\$s2}.
  \item \textbf{Instruction fields to Write Register~2}: The \texttt{rt} field
        (bits [20:16]) goes to Write Register~1, and the \texttt{rs} field
        (bits [25:21]) goes to Write Register~2 (or vice versa), so each port
        targets the correct register.
\end{enumerate}

\phantomsection\addcontentsline{toc}{subsubsection}{4.12.4}
\medskip\noindent\textbf{4.12.4}
The following new control signals are needed:
\begin{itemize}
  \item \textbf{Swap}: When asserted, enables both write ports and activates
        the cross-connected data routing (Read Data~1 $\to$ Write Data~2 and
        Read Data~2 $\to$ Write Data~1).
  \item \textbf{RegWrite2}: Enables the second write port. Can be tied to the
        Swap signal.
  \item \textbf{WriteDataSrc}: A mux select signal on the first write data
        port to choose between the normal write-back data (ALU result or
        memory data) and Read Data~2.
\end{itemize}

\phantomsection\addcontentsline{toc}{subsubsection}{4.12.5}
\medskip\noindent\textbf{4.12.5}
Modifications to Figure 4.21:
\begin{enumerate}
  \item Add a \textbf{second write port} to the Register File with inputs
        Write Register~2, Write Data~2, and RegWrite2.
  \item Add a \textbf{mux} before Write Data~1. Inputs: (0) normal write-back
        data (from the MemtoReg mux), (1) Read Data~2. Select signal: Swap.
  \item Connect \textbf{Read Data~1} directly to \textbf{Write Data~2}.
  \item Connect the \textbf{rs field} (instruction bits [25:21]) to Write
        Register~1 and the \textbf{rt field} (bits [20:16]) to Write
        Register~2 when Swap is asserted. This requires a mux on each Write
        Register input controlled by the Swap signal.
  \item The control unit decodes the \texttt{swap} opcode and asserts
        \textbf{Swap}, \textbf{RegWrite}, and \textbf{RegWrite2}, while
        deasserting MemRead, MemWrite, and Branch.
\end{enumerate}

%-----------------------------------------------------------------------------
\subsection*{4.13 --- Store Sum Instruction}
\hypertarget{sol-4.13}{}
\markright{4.13}
\addcontentsline{toc}{subsection}{4.13 --- Store Sum Instruction}

The proposed instruction is:
\begin{verbatim}
    ss $s1, $s2, 100($s3)  # Mem[$s3+100] = $s1 + $s2
\end{verbatim}

This instruction adds the contents of two registers and stores the result to a
memory address computed from a third register plus an offset.

\phantomsection\addcontentsline{toc}{subsubsection}{4.13.1}
\medskip\noindent\textbf{4.13.1}
We need \textbf{one new ALU (or adder)}. The existing ALU must compute the sum
\texttt{\$s1 + \$s2} (the data to be stored). Simultaneously, we need to
compute the memory address \texttt{\$s3 + sign\_ext(100)}. Since a single ALU
cannot perform both computations in the same cycle, a second ALU or dedicated
adder is required for the address calculation.

\phantomsection\addcontentsline{toc}{subsubsection}{4.13.2}
\medskip\noindent\textbf{4.13.2}
The \textbf{Register File} must be modified to support \textbf{three
simultaneous reads}. The standard MIPS Register File has two read ports, but
\texttt{ss} requires reading three registers: \texttt{\$s1}, \texttt{\$s2},
and \texttt{\$s3}. A third read port (Read Register~3, Read Data~3) must be
added.

Additionally, the instruction encoding must accommodate three register
specifiers and an immediate field, which is challenging in the standard 32-bit
MIPS format.

\phantomsection\addcontentsline{toc}{subsubsection}{4.13.3}
\medskip\noindent\textbf{4.13.3}
The following new data paths are needed:
\begin{enumerate}
  \item A path from the \textbf{third read port} (Read Data~3, carrying
        \texttt{\$s3}) to one input of the \textbf{new adder/ALU}.
  \item A path from the \textbf{sign-extended immediate} to the other input
        of the new adder/ALU (for computing the address).
  \item A path from the \textbf{existing ALU output} (\texttt{\$s1 + \$s2})
        to the \textbf{Write Data} input of Data Memory.
  \item A path from the \textbf{new adder/ALU output}
        (\texttt{\$s3 + sign\_ext(100)}) to the \textbf{Address} input of
        Data Memory.
\end{enumerate}

\phantomsection\addcontentsline{toc}{subsubsection}{4.13.4}
\medskip\noindent\textbf{4.13.4}
The following new control signals are needed:
\begin{itemize}
  \item \textbf{AddrSrc}: A mux select signal on the Data Memory address input
        to choose between the existing ALU output (for normal lw/sw) and the
        new adder output (for \texttt{ss}).
  \item \textbf{WriteDataSrc}: A mux select signal on the Data Memory write
        data input to choose between Read Data~2 (for normal \texttt{sw}) and
        the existing ALU output (for \texttt{ss}).
  \item \textbf{ALUSrc} must be deasserted so the existing ALU receives both
        register values (\texttt{\$s1} and \texttt{\$s2}) rather than a
        register and an immediate.
\end{itemize}

\phantomsection\addcontentsline{toc}{subsubsection}{4.13.5}
\medskip\noindent\textbf{4.13.5}
Modifications to Figure 4.21:
\begin{enumerate}
  \item Add a \textbf{third read port} to the Register File. The third
        register specifier comes from a new instruction field.
  \item Add a \textbf{second ALU/adder}. Its inputs are Read Data~3 and the
        sign-extended immediate. Its output is the memory address.
  \item Add a \textbf{mux before the Data Memory address} input:
        input~0 = existing ALU result (normal lw/sw), input~1 = new
        adder output (for \texttt{ss}). Select signal: AddrSrc.
  \item Add a \textbf{mux before the Data Memory write data} input:
        input~0 = Read Data~2 (normal \texttt{sw}),
        input~1 = existing ALU result (for \texttt{ss}).
        Select signal: WriteDataSrc.
  \item The control unit decodes the \texttt{ss} opcode and asserts
        MemWrite, AddrSrc, and WriteDataSrc, while deasserting RegWrite,
        MemRead, and Branch. ALUSrc~=~0 so the main ALU adds the two register
        values. ALUOp is set to ``add.''
\end{enumerate}

%-----------------------------------------------------------------------------
\subsection*{4.14 --- Imm Gen Block Critical Path}
\hypertarget{sol-4.14}{}
\markright{4.14}
\addcontentsline{toc}{subsection}{4.14 --- Imm Gen Block Critical Path}

We must determine whether the Imm Gen (sign-extend) block lies on the critical
path for any instruction. Recall the latencies:

\begin{center}
\begin{tabular}{ccccccccc}
\toprule
\textbf{I-Mem} & \textbf{Mux} & \textbf{ALU} & \textbf{Adder} &
\textbf{Gate} & \textbf{Reg} & \textbf{Reg File} &
\textbf{Sign ext.} & \textbf{Ctrl} \\
\midrule
200ps & 30ps & 120ps & 150ps & 5ps & 90ps & 150ps & 10ps & 50ps \\
\bottomrule
\end{tabular}
\end{center}

After the instruction is fetched from I-Mem (200\,ps), two operations begin in
parallel:
\begin{enumerate}
  \item \textbf{Register File read}: 150\,ps
  \item \textbf{Sign-extend}: 10\,ps (possibly also waiting for Control at
        50\,ps if the extension type depends on the opcode, but standard MIPS
        sign-extend is a fixed operation)
\end{enumerate}

The sign-extend completes in at most $10\,\text{ps}$ after I-Mem, while the
Register File read takes $150\,\text{ps}$. Both feed into the ALU input mux
(ALUSrc), so the mux cannot proceed until the \emph{slower} of the two
finishes. Since $150\,\text{ps} \gg 10\,\text{ps}$, the sign-extend result is
always ready long before the register read completes.

Let us verify by computing the critical paths for each instruction type:

\medskip\noindent\textbf{R-type:}
\[
  \underbrace{200}_{\text{I-Mem}} +
  \underbrace{150}_{\text{Reg File}} +
  \underbrace{30}_{\text{Mux}} +
  \underbrace{120}_{\text{ALU}} +
  \underbrace{30}_{\text{Mux}} +
  \underbrace{90}_{\text{Reg setup}}
  = 620\,\text{ps}
\]
Sign-extend is not used at all.

\medskip\noindent\textbf{Load (\texttt{lw}):}
\[
  \underbrace{200}_{\text{I-Mem}} +
  \underbrace{150}_{\text{Reg File}} +
  \underbrace{30}_{\text{Mux}} +
  \underbrace{120}_{\text{ALU}} +
  \underbrace{200}_{\text{D-Mem}} +
  \underbrace{30}_{\text{Mux}} +
  \underbrace{90}_{\text{Reg setup}}
  = 820\,\text{ps}
\]
Sign-extend feeds the ALUSrc mux, but it arrives at time $200 + 10 = 210$\,ps,
while the register value arrives at $200 + 150 = 350$\,ps. The mux waits for
the register value, so sign-extend is \emph{not} on the critical path.

\medskip\noindent\textbf{Store (\texttt{sw}):}
The address path is identical to \texttt{lw}; sign-extend is again dominated
by the register read.

\medskip\noindent\textbf{Branch (\texttt{beq}):}
Sign-extend feeds the branch target adder, arriving at $200 + 10 = 210$\,ps.
The PC+4 value from the other adder arrives at $200 + 150 = 350$\,ps (or the
register read finishes at 350\,ps for the comparator). Again, sign-extend is
not the bottleneck.

\medskip\noindent\textbf{Conclusion:} The Imm Gen (sign-extend) block is
\textbf{not on the critical path for any instruction}. Its 10\,ps latency is
always completely hidden by the much longer Register File read latency
(150\,ps) that operates in parallel on the same stage of the datapath.

%-----------------------------------------------------------------------------
\subsection*{4.15 --- No-Offset Load/Store Redesign}
\hypertarget{sol-4.15}{}
\markright{4.15}
\addcontentsline{toc}{subsection}{4.15 --- No-Offset Load/Store Redesign}

If \texttt{lw} and \texttt{sw} are modified to have no offset, the address is
taken directly from the \texttt{rs} register, and the ALU is bypassed for
memory instructions.

\phantomsection\addcontentsline{toc}{subsubsection}{4.15.1}
\medskip\noindent\textbf{4.15.1}
We compute the critical path for each instruction type under the new design.

\medskip\noindent\textit{R-type} (unchanged):
\[
  \underbrace{200}_{\text{I-Mem}} +
  \underbrace{150}_{\text{Reg File}} +
  \underbrace{30}_{\text{Mux}} +
  \underbrace{120}_{\text{ALU}} +
  \underbrace{30}_{\text{Mux}} +
  \underbrace{90}_{\text{Reg setup}}
  = 620\,\text{ps}
\]

\medskip\noindent\textit{Load (no-offset \texttt{lw}):}
The ALU is no longer on the load path. The register value goes directly to the
Data Memory address input:
\[
  \underbrace{200}_{\text{I-Mem}} +
  \underbrace{150}_{\text{Reg File}} +
  \underbrace{200}_{\text{D-Mem}} +
  \underbrace{30}_{\text{Mux}} +
  \underbrace{90}_{\text{Reg setup}}
  = 670\,\text{ps}
\]

\medskip\noindent\textit{Store (no-offset \texttt{sw}):}
\[
  \underbrace{200}_{\text{I-Mem}} +
  \underbrace{150}_{\text{Reg File}} +
  \underbrace{200}_{\text{D-Mem}}
  = 550\,\text{ps}
\]

\medskip\noindent\textit{Branch} (unchanged):
\[
  \underbrace{200}_{\text{I-Mem}} +
  \underbrace{150}_{\text{Reg File}} +
  \underbrace{30}_{\text{Mux}} +
  \underbrace{120}_{\text{ALU}} +
  \underbrace{5}_{\text{Gate}}
  = 505\,\text{ps}
\]

The new clock cycle time is determined by the longest path:
\[
  \text{New clock} = \max(620,\; 670,\; 550,\; 505) = \boxed{670\,\text{ps}}
\]

This is a reduction from the original 820\,ps (set by \texttt{lw} with offset).

\phantomsection\addcontentsline{toc}{subsubsection}{4.15.2}
\medskip\noindent\textbf{4.15.2}
Using the instruction mix from Exercise~4.3:

\begin{center}
\begin{tabular}{cccccc}
\toprule
\textbf{R-type} & \textbf{I-type (non-lw)} & \textbf{Load} &
\textbf{Store} & \textbf{Branch} & \textbf{Jump} \\
\midrule
24\% & 28\% & 25\% & 10\% & 11\% & 2\% \\
\bottomrule
\end{tabular}
\end{center}

On the \textbf{original CPU}: every instruction takes one cycle of 820\,ps.
For $N$ instructions:
\[
  T_{\text{old}} = N \times 820\,\text{ps}
\]

On the \textbf{new CPU}: each \texttt{lw} or \texttt{sw} that previously used
an offset must be replaced by two instructions (an \texttt{add}/\texttt{addi}
to compute the address, then the no-offset \texttt{lw}/\texttt{sw}). Load + Store
= $25\% + 10\% = 35\%$ of instructions. The new instruction count is:
\[
  N_{\text{new}} = N \times (1 + 0.35) = 1.35N
\]

Each instruction now takes one cycle of 670\,ps:
\[
  T_{\text{new}} = 1.35N \times 670\,\text{ps} = 904.5N\,\text{ps}
\]

Comparing:
\[
  \frac{T_{\text{new}}}{T_{\text{old}}}
  = \frac{904.5N}{820N}
  = 1.103
\]

The program runs approximately \textbf{10.3\% slower} on the new CPU.

\phantomsection\addcontentsline{toc}{subsubsection}{4.15.3}
\medskip\noindent\textbf{4.15.3}
The primary factor is the \textbf{fraction of \texttt{lw}/\texttt{sw}
instructions that actually use a non-zero offset}. If most memory instructions
already use the register directly (offset~=~0), then very few extra
instructions are needed and the faster clock cycle dominates, yielding a net
speedup. Conversely, if most memory instructions use non-zero offsets (as is
typical), the instruction count increase outweighs the clock improvement.

A secondary factor is the overall fraction of load/store instructions in the
program. Programs with few memory operations benefit more from the faster
clock because they add fewer extra instructions.

\phantomsection\addcontentsline{toc}{subsubsection}{4.15.4}
\medskip\noindent\textbf{4.15.4}
The \textbf{original CPU} (Figure 4.21) is generally the better overall design.
With a typical instruction mix (35\% loads and stores), the new design runs
about 10\% slower despite its faster clock. The offset field in \texttt{lw}/\texttt{sw}
is one of the most useful features of the MIPS ISA --- it allows direct access
to structure fields, stack variables, and array elements without separate
address arithmetic. Removing it increases code size, instruction count, and
register pressure, all of which degrade real-world performance.

The no-offset design would only be advantageous in unusual programs where
memory operations are very rare (less than about 18\% of instructions), making
the clock cycle reduction worth the extra instructions. For general-purpose
computing, the original design is superior.
%-----------------------------------------------------------------------------
\subsection*{4.16 --- Pipelining and Clock Cycle Time}
\hypertarget{sol-4.16}{}
\markright{4.16}
\addcontentsline{toc}{subsection}{4.16 --- Pipelining and Clock Cycle Time}

The pipeline stage latencies are:

\begin{center}
\begin{tabular}{ccccc}
\toprule
\textbf{IF} & \textbf{ID} & \textbf{EX} & \textbf{MEM} & \textbf{WB} \\
\midrule
250ps & 350ps & 150ps & 300ps & 200ps \\
\bottomrule
\end{tabular}
\end{center}

Instruction mix: ALU/Logic 40\%, Jump/Branch 20\%, Load 20\%, Store 20\%.

\phantomsection\addcontentsline{toc}{subsubsection}{4.16.1}
\medskip\noindent\textbf{4.16.1}
In a \textbf{pipelined} processor, every stage must fit within one clock cycle,
so the clock cycle time equals the latency of the slowest stage:
\[
  T_{\text{pipe}} = \max(250, 350, 150, 300, 200) = 350\text{ ps (ID stage).}
\]

In a \textbf{non-pipelined} (single-cycle) processor, the clock cycle must
accommodate every stage in sequence:
\[
  T_{\text{single}} = 250 + 350 + 150 + 300 + 200 = 1250\text{ ps.}
\]

\phantomsection\addcontentsline{toc}{subsubsection}{4.16.2}
\medskip\noindent\textbf{4.16.2}
The total latency of a single \texttt{lw} instruction (from fetch to
completion):

\begin{itemize}
  \item \textbf{Pipelined:} The instruction passes through all 5 stages, each
        taking the pipelined clock period:
        \[
          5 \times 350 = 1750\text{ ps.}
        \]
  \item \textbf{Non-pipelined:} The instruction takes exactly one (single-cycle)
        clock period:
        \[
          250 + 350 + 150 + 300 + 200 = 1250\text{ ps.}
        \]
\end{itemize}

\noindent Note: pipelining increases the latency of an individual instruction
but improves \emph{throughput} (one instruction completes per cycle once the
pipeline is full).

\phantomsection\addcontentsline{toc}{subsubsection}{4.16.3}
\medskip\noindent\textbf{4.16.3}
The clock cycle time is set by the slowest stage (ID at 350 ps). To reduce
it, we should split the ID stage into two stages, each with half the
latency:
\[
  \frac{350}{2} = 175\text{ ps per new stage.}
\]

The resulting 6-stage pipeline has latencies: 250, 175, 175, 150, 300, 200 ps.
The new clock cycle time is:
\[
  T_{\text{new}} = \max(250, 175, 175, 150, 300, 200) = 300\text{ ps (MEM stage).}
\]

\noindent Splitting any other stage would not help. For example, splitting
MEM ($300/2 = 150$ ps) gives latencies 250, 350, 150, 150, 150, 200, and
the clock is still limited by ID at 350 ps. Therefore, \textbf{split the ID
stage}; the new clock cycle time is \textbf{300 ps}.

\phantomsection\addcontentsline{toc}{subsubsection}{4.16.4}
\medskip\noindent\textbf{4.16.4}
The data memory is used during the MEM stage. Only \texttt{lw} and
\texttt{sw} instructions perform a useful data memory access. All other
instructions pass through MEM but do not read or write data memory.

\[
  \text{Data memory utilization} = \text{Load} + \text{Store}
    = 20\% + 20\% = \boxed{40\%.}
\]

\phantomsection\addcontentsline{toc}{subsubsection}{4.16.5}
\medskip\noindent\textbf{4.16.5}
The write-register port of the Registers unit is used during the WB stage.
Only instructions that write a result to a register use this port:
ALU/Logic instructions and Load instructions. Stores, branches, and jumps
do not write to a register.

\[
  \text{Write-register utilization} = \text{ALU/Logic} + \text{Load}
    = 40\% + 20\% = \boxed{60\%.}
\]

%-----------------------------------------------------------------------------
\subsection*{4.17 --- Minimum Cycles for N Instructions}
\hypertarget{sol-4.17}{}
\markright{4.17}
\addcontentsline{toc}{subsection}{4.17 --- Minimum Cycles for N Instructions}

A $k$-stage pipeline requires $k$ cycles to complete the first instruction.
After that, one instruction completes every cycle (assuming no stalls or
hazards). For $N$ instructions:
\[
  \text{Minimum cycles} = k + (N - 1).
\]

For a \textbf{4-stage pipeline}:
\[
  \text{Minimum cycles} = 4 + (N - 1) = N + 3.
\]

\noindent\textbf{Justification:} The first instruction enters the pipeline at
cycle 1 and completes at cycle 4. Each subsequent instruction completes one
cycle later, so the $N$-th instruction completes at cycle $4 + (N-1) = N+3$.

\begin{verbatim}
         CC1  CC2  CC3  CC4  CC5  CC6  ...  CC(N+3)
  I1:    S1   S2   S3   S4
  I2:         S1   S2   S3   S4
  I3:              S1   S2   S3   S4
  ...                             ...
  IN:                                       S4
\end{verbatim}

%-----------------------------------------------------------------------------
\subsection*{4.18 --- Pipeline without Hazard Detection}
\hypertarget{sol-4.18}{}
\markright{4.18}
\addcontentsline{toc}{subsection}{4.18 --- Pipeline without Hazard Detection}

Initial values: \texttt{\$s0} = 11, \texttt{\$s1} = 22.

\begin{verbatim}
    add $s0, $s0, $s1    # I1
    add $s2, $s0, $s1    # I2
\end{verbatim}

Pipeline execution (no hazard detection, no forwarding):

\begin{verbatim}
         CC1  CC2  CC3  CC4  CC5
  I1:    IF   ID   EX   MEM  WB
  I2:         IF   ID   EX   MEM  WB
\end{verbatim}

\noindent\textbf{Analysis:}

\begin{itemize}
  \item \textbf{I1} reads \texttt{\$s0} (= 11) and \texttt{\$s1} (= 22) in
        \textbf{CC2} (ID stage), computes $11 + 22 = 33$, and writes
        \texttt{\$s0} $\leftarrow$ 33 in \textbf{CC5} (WB stage).
  \item \textbf{I2} reads \texttt{\$s0} and \texttt{\$s1} in \textbf{CC3}
        (ID stage). At this point, I1 has not yet written back its result
        (that happens in CC5). So I2 reads the \emph{old} value of
        \texttt{\$s0} = 11.
  \item I2 computes $11 + 22 = 33$ and writes \texttt{\$s2} $\leftarrow$ 33
        in CC6.
\end{itemize}

\noindent\textbf{Final register values:}

\begin{center}
\begin{tabular}{lc}
\toprule
\textbf{Register} & \textbf{Value} \\
\midrule
\texttt{\$s0} & 33 \\
\texttt{\$s1} & 22 \\
\texttt{\$s2} & 33 (incorrect --- should be 55) \\
\bottomrule
\end{tabular}
\end{center}

\noindent Without hazard detection, I2 uses the stale value of \texttt{\$s0},
producing an incorrect result. The correct value of \texttt{\$s2} should be
$(11+22) + 22 = 55$.

%-----------------------------------------------------------------------------
\subsection*{4.19 --- Pipeline with Same-Cycle Forwarding}
\hypertarget{sol-4.19}{}
\markright{4.19}
\addcontentsline{toc}{subsection}{4.19 --- Pipeline with Same-Cycle Forwarding}

Initial values: \texttt{\$s0} = 11, \texttt{\$s1} = 22.

\begin{verbatim}
    add $s0, $s0, $s1    # I1
    nop                   # I2
    nop                   # I3
    add $s2, $s0, $s1    # I4
\end{verbatim}

Pipeline execution:

\begin{verbatim}
         CC1  CC2  CC3  CC4  CC5  CC6  CC7  CC8
  I1:    IF   ID   EX   MEM  WB
  nop:        IF   ID   EX   MEM  WB
  nop:             IF   ID   EX   MEM  WB
  I4:                   IF   ID   EX   MEM  WB
\end{verbatim}

\noindent\textbf{Analysis:}

\begin{itemize}
  \item \textbf{I1} computes \texttt{\$s0} = $11 + 22 = 33$ and writes back
        in \textbf{CC5} (WB stage).
  \item \textbf{I4} reads registers in \textbf{CC5} (ID stage).
  \item With \emph{same-cycle forwarding} (also called WB$\to$ID forwarding),
        the register file is designed so that a write in the first half of the
        cycle is available to a read in the second half of the same cycle.
        Therefore, I4 reads the \emph{new} value \texttt{\$s0} = 33 in CC5.
  \item I4 computes \texttt{\$s2} = $33 + 22 = 55$.
\end{itemize}

\noindent\textbf{Final register values:}

\begin{center}
\begin{tabular}{lc}
\toprule
\textbf{Register} & \textbf{Value} \\
\midrule
\texttt{\$s0} & 33 \\
\texttt{\$s1} & 22 \\
\texttt{\$s2} & 55 (correct) \\
\bottomrule
\end{tabular}
\end{center}

\noindent The two \texttt{nop} instructions ensure that I1's WB stage and I4's
ID stage occur in the same cycle, allowing the forwarded value to be used.

%-----------------------------------------------------------------------------
\subsection*{4.20 --- Adding NOPs for Data Hazards}
\hypertarget{sol-4.20}{}
\markright{4.20}
\addcontentsline{toc}{subsection}{4.20 --- Adding NOPs for Data Hazards}

Original code (no hazard handling, no forwarding):

\begin{verbatim}
    add $s0, $s1, $s5    # I1: writes $s0
    add $s2, $s0, $s5    # I2: reads $s0 (from I1)
    add $s4, $s0, $s1    # I3: reads $s0 (from I1)
\end{verbatim}

\noindent\textbf{Dependency analysis:}
\begin{itemize}
  \item I2 depends on I1 via \texttt{\$s0} (RAW hazard, distance = 1 instruction).
  \item I3 depends on I1 via \texttt{\$s0} (RAW hazard, distance = 2 instructions).
  \item I3 does \emph{not} depend on I2 (I3 reads \texttt{\$s0} and
        \texttt{\$s1}; I2 writes \texttt{\$s2}).
\end{itemize}

Without forwarding, a result written in the WB stage (cycle $i + 4$) is
available to a subsequent instruction's ID stage only if there are at least 2
instructions (or \texttt{nop}s) between the producer and consumer (so that WB
and ID overlap in the same cycle, assuming same-cycle forwarding in the
register file).

We need 2 \texttt{nop}s between I1 and I2:

\begin{verbatim}
    add $s0, $s1, $s5    # I1: writes $s0
    nop
    nop
    add $s2, $s0, $s5    # I2: reads $s0 (now safe)
    add $s4, $s0, $s1    # I3: reads $s0 (safe -- 4 instr. after I1)
\end{verbatim}

\noindent\textbf{Verification with pipeline diagram:}

\begin{verbatim}
         CC1  CC2  CC3  CC4  CC5  CC6  CC7  CC8  CC9
  I1:    IF   ID   EX   MEM  WB
  nop:        IF   ID   EX   MEM  WB
  nop:             IF   ID   EX   MEM  WB
  I2:                   IF   ID   EX   MEM  WB
  I3:                        IF   ID   EX   MEM  WB
\end{verbatim}

\begin{itemize}
  \item I1 writes \texttt{\$s0} in CC5 (WB). I2 reads \texttt{\$s0} in CC5
        (ID). With same-cycle register file forwarding, I2 gets the correct
        value. \checkmark
  \item I3 reads \texttt{\$s0} in CC6 (ID), well after I1 writes in CC5.
        \checkmark
  \item No dependency exists between I2 and I3, so no additional
        \texttt{nop}s are needed.
\end{itemize}

\noindent\textbf{Answer:} Insert 2 \texttt{nop}s between the first and second
instructions. No \texttt{nop}s are needed between the second and third
instructions.

%-----------------------------------------------------------------------------
\subsection*{4.21 --- Forwarding Hardware Speedup}
\hypertarget{sol-4.21}{}
\markright{4.21}
\addcontentsline{toc}{subsection}{4.21 --- Forwarding Hardware Speedup}

\noindent\textbf{Given:}
\begin{itemize}
  \item Without forwarding: clock cycle = 250 ps, 5\% of instructions are
        \texttt{nop}s inserted to resolve hazards.
  \item With forwarding: clock cycle = 300 ps (increased due to forwarding
        hardware), only 2\% \texttt{nop}s needed.
\end{itemize}

\medskip\noindent\textbf{Execution time comparison:}

For a program that needs $N$ ``useful'' instructions, the total instructions
executed (including \texttt{nop}s) and execution times are:

\begin{align*}
  T_{\text{no fwd}} &= 250 \times N \times 1.05 = 262.5\,N \text{ ps,} \\
  T_{\text{fwd}}    &= 300 \times N \times 1.02 = 306.0\,N \text{ ps.}
\end{align*}

\noindent\textbf{Speedup} of the forwarding pipeline over the non-forwarding
pipeline:

\[
  \text{Speedup} = \frac{T_{\text{no fwd}}}{T_{\text{fwd}}}
    = \frac{262.5}{306.0}
    \approx 0.858.
\]

Since the speedup is less than 1, the forwarding hardware actually makes the
processor \textbf{slower}, not faster. The 20\% increase in clock cycle time
(from 250 to 300 ps) outweighs the reduction in \texttt{nop}s (from 5\% to
2\%).

\medskip\noindent\textbf{Conclusion:} The forwarding hardware results in a
\textbf{slowdown of approximately 14.2\%}. In this scenario, adding forwarding
is not worthwhile because the clock penalty exceeds the benefit of fewer stalls.
%-----------------------------------------------------------------------------
\subsection*{4.22 --- Structural Hazard: Unified Memory}
\hypertarget{sol-4.22}{}
\markright{4.22}
\addcontentsline{toc}{subsection}{4.22 --- Structural Hazard: Unified Memory}

With a unified memory, the instruction fetch (IF) stage and the data memory
access (MEM) stage share a single memory port.  A \emph{structural hazard}
arises whenever one instruction is in IF while another instruction is in MEM
during the same clock cycle.

\phantomsection\addcontentsline{toc}{subsubsection}{4.22.1}
\medskip\noindent\textbf{4.22.1}
In the standard 5-stage pipeline the MEM stage of an instruction occurs three
cycles after its IF stage.  Therefore, the structural hazard occurs between
an instruction in MEM and the instruction that enters IF three cycles later.

For the given three-instruction sequence, consider what happens when
additional instructions follow:

\begin{verbatim}
             CC1  CC2  CC3  CC4  CC5  CC6  CC7  CC8  CC9
    lw  $s5:  IF   ID   EX  MEM   WB
    sw  $s4:       IF   ID   EX  MEM   WB
    add $s2:            IF   ID   EX  MEM   WB
    nxt1:           stall--> IF   ID   EX  MEM   WB
    nxt2:                stall--> IF   ID   EX  MEM  WB
\end{verbatim}

\noindent Explanation:
\begin{itemize}
  \item \textbf{CC4}: \texttt{lw} is in MEM (reading data memory).  The 4th
        instruction (\texttt{nxt1}) would normally begin IF in CC4, but the
        memory is busy $\Rightarrow$ \textbf{stall}.  \texttt{nxt1} does IF in
        CC5 instead.
  \item \textbf{CC5}: \texttt{sw} is in MEM (writing data memory).  The 5th
        instruction (\texttt{nxt2}) would normally begin IF in CC5, but the
        memory is again busy $\Rightarrow$ \textbf{stall}.  \texttt{nxt2} does
        IF in CC6 instead.
  \item \texttt{add} is in MEM in CC6, but \texttt{add} is an R-type
        instruction: it does \emph{not} access data memory, so the memory port
        is free for a subsequent instruction's IF.  No stall.
\end{itemize}

Among the three given instructions alone, no structural hazard manifests
(their IF and MEM stages do not overlap with each other).  The stalls affect
the instructions that \emph{follow} the loads and stores.

\phantomsection\addcontentsline{toc}{subsubsection}{4.22.2}
\medskip\noindent\textbf{4.22.2}
Yes, reordering can reduce (but not always eliminate) these stalls.  The key
insight is that only \texttt{lw} and \texttt{sw} instructions actually use
data memory in the MEM stage; R-type and branch instructions do not.  If we
schedule non-memory instructions so that their MEM stages coincide with the
IF of later instructions, the structural hazard is avoided.

For example, reordering the code so that a non-memory instruction separates
two consecutive memory instructions can prevent their MEM stages from
blocking fetches of nearby instructions.

\phantomsection\addcontentsline{toc}{subsubsection}{4.22.3}
\medskip\noindent\textbf{4.22.3}
This structural hazard \textbf{must} be handled in hardware (via stall logic
or by providing separate instruction and data memories/caches).  It
\textbf{cannot} be eliminated purely by adding \texttt{nop} instructions,
because:
\begin{itemize}
  \item Every \texttt{nop} must itself be fetched from instruction memory.
  \item A \texttt{nop} in the MEM stage does not use data memory, so it does
        not cause a conflict.  However, the \texttt{nop} still needs to be
        fetched (IF), and that fetch may conflict with the MEM of another
        instruction.
  \item The fundamental issue is that with a single memory port, the hardware
        must arbitrate between IF and MEM.  Software scheduling can reduce the
        frequency of conflicts, but there is no software-only guarantee that a
        \texttt{nop}'s own IF will not collide with another instruction's MEM.
\end{itemize}

In practice, the simplest hardware solution is to stall the pipeline (delay
the IF stage) whenever a MEM access is occurring.  The best long-term
solution is to use separate instruction and data caches (a Harvard-style
memory architecture), which eliminates the hazard entirely.

\phantomsection\addcontentsline{toc}{subsubsection}{4.22.4}
\medskip\noindent\textbf{4.22.4}
A structural hazard stall occurs whenever an instruction's MEM stage actually
accesses data memory, which happens only for \texttt{lw} and \texttt{sw}
instructions.  From the instruction mix in Exercise~4.3:
\[
  \text{Fraction using data memory} = \text{Load} + \text{Store}
    = 25\% + 10\% = 35\%.
\]
Each such instruction causes one stall cycle (the IF of the instruction
three positions later is blocked for one cycle).  Therefore, approximately
\textbf{35\% of instructions will cause a one-cycle stall}, adding roughly
0.35 extra cycles per instruction to the CPI.

%-----------------------------------------------------------------------------
\subsection*{4.23 --- Four-Stage Pipeline}
\hypertarget{sol-4.23}{}
\markright{4.23}
\addcontentsline{toc}{subsection}{4.23 --- Four-Stage Pipeline}

If load/store instructions use a register directly as the memory address
(with no offset), those instructions no longer require the ALU to compute
an effective address.  Consequently, the EX (ALU) and MEM (data memory)
stages can be merged into a single stage, yielding a 4-stage pipeline:
\[
  \text{IF} \;\to\; \text{ID} \;\to\; \text{EX/MEM} \;\to\; \text{WB}
\]

\phantomsection\addcontentsline{toc}{subsubsection}{4.23.1}
\medskip\noindent\textbf{4.23.1}
The reduction from 5 stages to 4 stages affects performance in several ways:
\begin{itemize}
  \item The \textbf{pipeline latency} for each instruction drops from 5 cycles
        to 4 cycles.  This reduces startup cost and branch penalty.
  \item The \textbf{ideal CPI} is still 1 (one instruction completes per
        cycle in steady state).
  \item \textbf{Data hazard penalties decrease.}  In a 5-stage pipeline, a
        load-use hazard requires 1 stall (the load result is available after
        MEM, but the dependent instruction needs it in EX).  In the 4-stage
        pipeline, the load produces its result at the end of the combined
        EX/MEM stage, and the dependent instruction needs it at the start of
        the \emph{next} EX/MEM stage---so forwarding from the EX/MEM output to
        the EX/MEM input works with no stall.  This eliminates load-use
        hazards entirely.
  \item \textbf{Branch penalty decreases.}  With branches resolved in the
        EX/MEM stage (stage 3), the misprediction penalty is 2 cycles instead
        of 3 (in a 5-stage pipeline where branches resolve in EX, the penalty
        is 2; if resolved in MEM, it is 3).  In the 4-stage pipeline the
        penalty is at most 2 cycles (stages IF and ID must be flushed).
  \item \textbf{Fewer pipeline registers} are needed (3 sets instead of 4),
        reducing hardware cost and power.
\end{itemize}

\phantomsection\addcontentsline{toc}{subsubsection}{4.23.2}
\medskip\noindent\textbf{4.23.2}
This change can \textbf{improve} pipeline performance because:
\begin{enumerate}
  \item \textbf{Elimination of load-use stalls:} As noted above, forwarding
        from the end of the combined EX/MEM stage to its own input removes
        the 1-cycle load-use penalty that plagues the 5-stage design.  Since
        loads are $\approx 25\%$ of instructions and many are followed by
        dependent instructions, this can significantly reduce CPI.
  \item \textbf{Reduced branch penalty:} Fewer stages between IF and the
        branch resolution stage means fewer wasted cycles on mispredictions.
  \item \textbf{Simpler forwarding logic:} Fewer stages mean fewer forwarding
        paths and a simpler hazard detection unit.
\end{enumerate}

\phantomsection\addcontentsline{toc}{subsubsection}{4.23.3}
\medskip\noindent\textbf{4.23.3}
This change can \textbf{degrade} performance because:
\begin{enumerate}
  \item \textbf{Longer clock period:} The combined EX/MEM stage must
        accommodate the latency of both the ALU \emph{and} the data memory
        access (for load/store) or just the ALU (for R-type).  The clock
        period is set by the slowest stage, so:
        \[
          T_{\text{clock}} = \max(\text{IF},\; \text{ID},\;
            \text{ALU} + \text{D-Mem},\; \text{WB}).
        \]
        Using the latencies from Exercise~4.7:
        $\text{ALU} + \text{D-Mem} = 120 + 200 = 320\;\text{ps}$ (at
        minimum), which is likely the longest stage and thus increases the
        clock period.
  \item \textbf{Increased instruction count:} Without address offsets,
        programs that currently use \texttt{lw \$t0, 8(\$s1)} must first
        compute \texttt{\$s1 + 8} in a separate \texttt{addi} instruction,
        then perform the load.  This increases the dynamic instruction count.
        From Exercise~4.3, loads and stores are $25\% + 10\% = 35\%$ of
        instructions; many of these use non-zero offsets, so the instruction
        count could increase substantially.
  \item \textbf{Reduced throughput:} The longer clock period directly reduces
        the instruction throughput (instructions per second), partially or
        fully offsetting the gains from fewer hazards.
\end{enumerate}

%-----------------------------------------------------------------------------
\subsection*{4.24 --- Pipeline Diagram Selection}
\hypertarget{sol-4.24}{}
\markright{4.24}
\addcontentsline{toc}{subsection}{4.24 --- Pipeline Diagram Selection}

The two pipeline diagrams are:

\begin{verbatim}
    (a) ld  x1, x12:     IF  ID  EX  ME  WB
        add x5, x1, x17:     IF  ID  EX  ME  WB

    (b) ld  x1, x12:     IF  ID  EX  ME  WB
        add x5, x1, x17:         IF  ID  EX  ME  WB
\end{verbatim}

\textbf{Diagram (b)} correctly represents the operation of the hazard
detection unit.  The reasoning is as follows:

\begin{enumerate}
  \item The \texttt{add} instruction reads register \texttt{x1}, which is the
        destination of the preceding \texttt{ld} instruction.  This is a
        \emph{load-use} (RAW) data hazard.

  \item In the standard 5-stage pipeline with forwarding, the \texttt{ld}
        instruction does not produce the loaded value until the end of its
        MEM stage.  With forwarding, this value can be sent from the MEM/WB
        pipeline register to the EX stage of the next instruction---but only
        one cycle \emph{after} the load's MEM stage.

  \item In diagram~(a), \texttt{add} enters ID in CC3 and EX in CC4.
        However, \texttt{ld} is in EX during CC3 and does not reach MEM until
        CC4.  The loaded value is not available until the \emph{end} of CC4,
        but \texttt{add} needs it at the \emph{beginning} of CC4 (start of
        its EX stage).  Even with forwarding, the value is not ready in time.

  \item The hazard detection unit recognizes this conflict during CC3
        (when \texttt{ld} is in EX and \texttt{add} is in ID) and inserts a
        \textbf{one-cycle stall} (pipeline bubble).

  \item In diagram~(b), the stall shifts \texttt{add} by one cycle:
        \texttt{add} does IF in CC2, then is stalled in CC3 (bubble inserted),
        repeats ID in CC4, and enters EX in CC5.  Now \texttt{ld} finishes
        MEM at the end of CC4, and the MEM/WB forwarding path delivers the
        loaded value to \texttt{add}'s EX stage at the start of CC5.
        \checkmark
\end{enumerate}

Therefore, diagram~(b) is correct because the hazard detection unit must
insert a one-cycle stall for any load-use hazard, even when full forwarding
is available.

%-----------------------------------------------------------------------------
\subsection*{4.25 --- Loop Pipeline Execution}
\hypertarget{sol-4.25}{}
\markright{4.25}
\addcontentsline{toc}{subsection}{4.25 --- Loop Pipeline Execution}

\phantomsection\addcontentsline{toc}{subsubsection}{4.25.1}
\medskip\noindent\textbf{4.25.1}
The loop body is:
\begin{verbatim}
    LOOP: lw   $s0, 0($s3)       # I1
          add  $s1, $s0, $s2     # I2  (depends on I1: $s0)
          sw   $s1, 0($s3)       # I3  (depends on I2: $s1)
          addi $s3, $s3, -4      # I4
          bne  $s3, $zero, LOOP  # I5  (depends on I4: $s3)
\end{verbatim}

\noindent \textbf{Assumptions:}
\begin{itemize}
  \item Full forwarding (EX/MEM $\to$ EX and MEM/WB $\to$ EX).
  \item Branches resolved in the EX stage; perfect branch prediction (no
        control hazard stalls).
  \item No delay slots.
\end{itemize}

\noindent \textbf{Dependency analysis:}
\begin{itemize}
  \item \texttt{lw} $\to$ \texttt{add} (\texttt{\$s0}):  \textbf{Load-use
        hazard}.  \texttt{lw} produces \texttt{\$s0} at end of MEM;
        \texttt{add} needs it at start of EX.  Even with MEM/WB forwarding,
        there is a 1-cycle stall.
  \item \texttt{add} $\to$ \texttt{sw} (\texttt{\$s1}):  EX/MEM forwarding.
        \texttt{add} produces \texttt{\$s1} at end of EX; \texttt{sw} reads
        \texttt{\$s1} in EX (for the store data path).  Forwarding from
        EX/MEM register works with no stall.
  \item \texttt{addi} $\to$ \texttt{bne} (\texttt{\$s3}):  \texttt{addi}
        produces \texttt{\$s3} at end of EX.  \texttt{bne} needs \texttt{\$s3}
        at the start of its EX stage.  \texttt{addi} is immediately before
        \texttt{bne}, so EX/MEM forwarding delivers the result with no stall.
  \item \texttt{bne} $\to$ next iteration's \texttt{lw} (\texttt{\$s3}):
        No dependency on \texttt{\$s3} via forwarding is needed because
        \texttt{addi} wrote \texttt{\$s3} and it will be in the register file
        or forwarded by the time the next \texttt{lw} reaches EX.
\end{itemize}

\noindent \textbf{Pipeline execution diagram} (two iterations):

{\small
\begin{verbatim}
Cycle:       1    2    3    4    5    6    7    8    9   10   11   12   13   14
lw   $s0:   IF   ID   EX  MEM   WB
add  $s1:        IF   ID   **   EX  MEM   WB
sw   $s1:             IF   **   ID   EX  MEM   WB
addi $s3:                  IF   ID   EX  MEM   WB
bne:                            IF   ID   EX  MEM   WB
lw'  $s0:                           IF   ID   EX  MEM   WB
add' $s1:                                IF   ID   **   EX  MEM   WB
sw'  $s1:                                     IF   **   ID   EX  MEM   WB
addi'$s3:                                          IF   ID   EX  MEM   WB
bne':                                                   IF   ID   EX  MEM  WB
\end{verbatim}
}

\noindent where \texttt{**} denotes a one-cycle stall (bubble) due to the
load-use hazard between \texttt{lw} and \texttt{add}.

\noindent \textbf{Cycle-by-cycle explanation (first iteration):}
\begin{itemize}
  \item \textbf{CC1}: \texttt{lw} -- IF.
  \item \textbf{CC2}: \texttt{lw} -- ID; \texttt{add} -- IF.
  \item \textbf{CC3}: \texttt{lw} -- EX; \texttt{add} -- ID (hazard
        detected: \texttt{lw} in EX is a load, \texttt{add} in ID reads
        \texttt{\$s0}).  Pipeline stalls \texttt{add} and all instructions
        behind it.  \texttt{sw} -- IF.
  \item \textbf{CC4}: \texttt{lw} -- MEM; bubble in EX; \texttt{add} -- ID
        (held); \texttt{sw} -- IF (held); \texttt{addi} would enter IF but is
        stalled.
  \item \textbf{CC5}: \texttt{lw} -- WB; \texttt{add} -- EX (receives
        \texttt{\$s0} via MEM$\to$EX forwarding); \texttt{sw} -- ID;
        \texttt{addi} -- IF.
  \item \textbf{CC6}: \texttt{add} -- MEM; \texttt{sw} -- EX (receives
        \texttt{\$s1} via EX$\to$EX forwarding); \texttt{addi} -- ID;
        \texttt{bne} -- IF.
  \item \textbf{CC7}: \texttt{add} -- WB; \texttt{sw} -- MEM;
        \texttt{addi} -- EX; \texttt{bne} -- ID; next iteration
        \texttt{lw'} -- IF (branch predicted taken).
  \item \textbf{CC8}: \texttt{sw} -- WB; \texttt{addi} -- MEM;
        \texttt{bne} -- EX (receives \texttt{\$s3} from \texttt{addi} via
        EX$\to$EX forwarding; branch resolved, prediction correct);
        \texttt{lw'} -- ID.
  \item \textbf{CC9}: \texttt{addi} -- WB; \texttt{bne} -- MEM;
        \texttt{lw'} -- EX; \texttt{add'} -- ID (hazard detected again).
\end{itemize}

\noindent \textbf{Summary:} Each iteration takes \textbf{6 cycles}
(5~instructions $+$ 1~stall from the load-use hazard).  For $n$ iterations
the total cycle count is approximately $6n + 4$ (6 cycles per iteration
plus 4 cycles of pipeline startup/drain).
%-----------------------------------------------------------------------------
\subsection*{4.26 --- RAW Dependencies and Forwarding}
\hypertarget{sol-4.26}{}
\markright{4.26}
\addcontentsline{toc}{subsection}{4.26 --- RAW Dependencies and Forwarding}

This problem uses the pipelined datapath from Figure~4.33, with the following
RAW data-dependence frequencies:

\begin{center}
\begin{tabular}{cccccc}
\toprule
\textbf{EX to 1st} & \textbf{MEM to 1st} & \textbf{EX to 2nd} &
\textbf{MEM to 2nd} & \textbf{EX to 1st \&} & \textbf{MEM to 1st \&} \\
& & & & \textbf{EX to 2nd} & \textbf{MEM to 2nd} \\
\midrule
20\% & 5\% & 20\% & 10\% & 10\% & 5\% \\
\bottomrule
\end{tabular}
\end{center}

\phantomsection\addcontentsline{toc}{subsubsection}{4.26.1}
\medskip\noindent\textbf{4.26.1}
Example instruction sequences for each RAW dependency type:

\textit{EX to 1st} --- the producing instruction's ALU result is needed by the
very next instruction:
\begin{verbatim}
    add $s0, $s1, $s2     # produces $s0 in EX
    add $s3, $s0, $s4     # consumes $s0 in EX (1st after producer)
    add $s5, $s6, $s7
\end{verbatim}

\textit{MEM to 1st} --- a load's memory result is needed by the very next
instruction:
\begin{verbatim}
    lw  $s0, 0($s1)       # produces $s0 in MEM
    add $s3, $s0, $s4     # consumes $s0 in EX (1st after producer)
    add $s5, $s6, $s7
\end{verbatim}

\textit{EX to 2nd} --- the producing instruction's ALU result is needed by the
instruction two later:
\begin{verbatim}
    add $s0, $s1, $s2     # produces $s0 in EX
    add $s5, $s6, $s7     # no dependency
    add $s3, $s0, $s4     # consumes $s0 in EX (2nd after producer)
\end{verbatim}

\textit{MEM to 2nd} --- a load's memory result is needed by the instruction two
later:
\begin{verbatim}
    lw  $s0, 0($s1)       # produces $s0 in MEM
    add $s5, $s6, $s7     # no dependency
    add $s3, $s0, $s4     # consumes $s0 in EX (2nd after producer)
\end{verbatim}

\textit{EX to 1st and EX to 2nd} --- the same ALU result is needed by both the
next and the second-next instruction:
\begin{verbatim}
    add $s0, $s1, $s2     # produces $s0 in EX
    add $s3, $s0, $s4     # consumes $s0 (1st after producer)
    add $s5, $s0, $s6     # consumes $s0 (2nd after producer)
\end{verbatim}

\textit{MEM to 1st and MEM to 2nd} --- a load result is needed by both the next
and the second-next instruction:
\begin{verbatim}
    lw  $s0, 0($s1)       # produces $s0 in MEM
    add $s3, $s0, $s4     # consumes $s0 (1st after producer)
    add $s5, $s0, $s6     # consumes $s0 (2nd after producer)
\end{verbatim}

\phantomsection\addcontentsline{toc}{subsubsection}{4.26.2}
\medskip\noindent\textbf{4.26.2}
Number of \texttt{nop}s needed for each dependency type on a pipeline with
\emph{no} forwarding and \emph{no} hazard detection (the value must be written
back to the register file before the consumer reads it in ID; we assume the
register file supports write-first-half / read-second-half in the same cycle):

\begin{center}
\begin{tabular}{lc}
\toprule
\textbf{Dependency type} & \textbf{NOPs needed} \\
\midrule
EX to 1st              & 2 \\
MEM to 1st             & 2 \\
EX to 2nd              & 1 \\
MEM to 2nd             & 1 \\
EX to 1st \& EX to 2nd   & 2 \\
MEM to 1st \& MEM to 2nd & 2 \\
\bottomrule
\end{tabular}
\end{center}

\textit{Reasoning:} Without forwarding, the consumer must wait until the producer
has completed WB (or at least the first half of the WB cycle) so that the value
is available in the register file for the consumer's ID stage. For both
``to 1st'' types the consumer is one slot behind the producer, so we insert 2
NOPs to create a 2-instruction gap (the producer reaches WB while the consumer
is in ID). For both ``to 2nd'' types there is already a 1-instruction gap, so
only 1 NOP is needed. For the combined cases, the tightest constraint is the
``to 1st'' dependency, which requires 2 NOPs; those same NOPs also satisfy the
``to 2nd'' constraint.

\phantomsection\addcontentsline{toc}{subsubsection}{4.26.3}
\medskip\noindent\textbf{4.26.3}
Consider the following three-instruction sequence:
\begin{verbatim}
    add $s0, $s1, $s2     # I1: produces $s0
    add $s3, $s0, $s4     # I2: reads $s0 (EX to 1st from I1)
    add $s5, $s0, $s6     # I3: reads $s0 (EX to 2nd from I1)
\end{verbatim}

Analyzing each dependency \emph{independently}:
\begin{itemize}
  \item I1 $\to$ I2 (EX to 1st): 2 stalls if considered alone.
  \item I1 $\to$ I3 (EX to 2nd): 1 stall if considered alone.
  \item Sum of individual stalls $= 2 + 1 = 3$.
\end{itemize}

However, the 2 NOPs inserted after I1 to resolve the I1$\to$I2 dependency also
push I3 further from I1 (I3 is now 4 instructions after I1), which more than
satisfies the I1$\to$I3 requirement. Therefore the actual stalls needed $= 2$,
which is less than the sum of $3$.

\phantomsection\addcontentsline{toc}{subsubsection}{4.26.4}
\medskip\noindent\textbf{4.26.4}
We analyze each forwarding configuration separately. The table frequencies use
inclusive counting (the ``combined'' entries are subsets of the individual
entries), so we decompose into exclusive categories:

\begin{center}
\begin{tabular}{lc}
\toprule
\textbf{Exclusive category} & \textbf{Frequency} \\
\midrule
EX to 1st only (not also EX to 2nd)       & $20\% - 10\% = 10\%$ \\
EX to 2nd only (not also EX to 1st)       & $20\% - 10\% = 10\%$ \\
EX to 1st \& EX to 2nd                    & $10\%$ \\
MEM to 1st only (not also MEM to 2nd)     & $5\% - 5\% = 0\%$ \\
MEM to 2nd only (not also MEM to 1st)     & $10\% - 5\% = 5\%$ \\
MEM to 1st \& MEM to 2nd                  & $5\%$ \\
\bottomrule
\end{tabular}
\end{center}

\textbf{No forwarding baseline.}
Stall cycles per category (using the results from 4.26.2):

\begin{center}
\begin{tabular}{lcc}
\toprule
\textbf{Category} & \textbf{Stalls} & \textbf{Contribution} \\
\midrule
EX to 1st only     & 2 & $0.10 \times 2 = 0.20$ \\
EX to 2nd only     & 1 & $0.10 \times 1 = 0.10$ \\
EX to 1st \& 2nd   & 2 & $0.10 \times 2 = 0.20$ \\
MEM to 1st only    & 2 & $0.00 \times 2 = 0.00$ \\
MEM to 2nd only    & 1 & $0.05 \times 1 = 0.05$ \\
MEM to 1st \& 2nd  & 2 & $0.05 \times 2 = 0.10$ \\
\midrule
\multicolumn{2}{l}{Total extra CPI} & $0.65$ \\
\bottomrule
\end{tabular}
\end{center}
\[
  \text{CPI}_{\text{no fwd}} = 1 + 0.65 = 1.65
\]

\textbf{EX-EX forwarding only}
(forward from EX/MEM pipeline register to ALU inputs in EX stage).

This eliminates stalls for any dependency where the result is produced in EX
and the consumer is the very next instruction (EX to 1st). For EX to 2nd,
by the time the consumer reaches EX, the producer is in WB; the value is
written to the register file in the first half of WB, and the consumer reads
it in the second half of ID (one stage earlier). With the write-first /
read-second convention, EX to 2nd is resolved through the register file
without needing MEM-EX forwarding. For MEM (load) dependencies, EX-EX
forwarding does not help since the data is not available until after MEM.

\begin{center}
\begin{tabular}{lcc}
\toprule
\textbf{Category} & \textbf{Stalls} & \textbf{Contribution} \\
\midrule
EX to 1st only     & 0 & $0.10 \times 0 = 0.00$ \\
EX to 2nd only     & 0 & $0.10 \times 0 = 0.00$ \\
EX to 1st \& 2nd   & 0 & $0.10 \times 0 = 0.00$ \\
MEM to 1st only    & 2 & $0.00 \times 2 = 0.00$ \\
MEM to 2nd only    & 1 & $0.05 \times 1 = 0.05$ \\
MEM to 1st \& 2nd  & 2 & $0.05 \times 2 = 0.10$ \\
\midrule
\multicolumn{2}{l}{Total extra CPI} & $0.15$ \\
\bottomrule
\end{tabular}
\end{center}
\[
  \text{CPI}_{\text{EX-EX only}} = 1 + 0.15 = 1.15
\]

\textbf{MEM-EX forwarding only}
(forward from MEM/WB pipeline register to ALU inputs in EX stage).

This eliminates stalls where the consumer can get the value from MEM/WB
(MEM to 1st after a 1-cycle stall is actually still needed since we don't
have EX-EX forwarding). With only MEM-EX forwarding:
\begin{itemize}
  \item EX to 1st: no EX-EX forwarding available. Need 1 stall so producer
    moves from EX/MEM to MEM/WB, then forward via MEM-EX path.
  \item EX to 2nd: producer is in MEM/WB when consumer reaches EX $\to$
    forwarded via MEM-EX path $\to$ 0 stalls.
  \item MEM to 1st: need 1 stall (producer must complete MEM before data is
    in MEM/WB; after 1 stall, MEM-EX forwarding delivers the load value) $\to$
    1 stall.
  \item MEM to 2nd: producer is in WB when consumer is in EX; value already
    written to register file $\to$ 0 stalls.
\end{itemize}

\begin{center}
\begin{tabular}{lcc}
\toprule
\textbf{Category} & \textbf{Stalls} & \textbf{Contribution} \\
\midrule
EX to 1st only     & 1 & $0.10 \times 1 = 0.10$ \\
EX to 2nd only     & 0 & $0.10 \times 0 = 0.00$ \\
EX to 1st \& 2nd   & 1 & $0.10 \times 1 = 0.10$ \\
MEM to 1st only    & 1 & $0.00 \times 1 = 0.00$ \\
MEM to 2nd only    & 0 & $0.05 \times 0 = 0.00$ \\
MEM to 1st \& 2nd  & 1 & $0.05 \times 1 = 0.05$ \\
\midrule
\multicolumn{2}{l}{Total extra CPI} & $0.25$ \\
\bottomrule
\end{tabular}
\end{center}
\[
  \text{CPI}_{\text{MEM-EX only}} = 1 + 0.25 = 1.25
\]

\phantomsection\addcontentsline{toc}{subsubsection}{4.26.5}
\medskip\noindent\textbf{4.26.5}
With full forwarding (both EX-EX and MEM-EX paths), all RAW hazards are
resolved except load-use hazards (MEM to 1st), which still require 1 stall
cycle because the load data is not available until the end of the MEM stage,
one cycle too late for the immediately following instruction's EX stage.

\begin{center}
\begin{tabular}{lcc}
\toprule
\textbf{Category} & \textbf{Stalls} & \textbf{Contribution} \\
\midrule
EX to 1st only     & 0 & $0.10 \times 0 = 0.00$ \\
EX to 2nd only     & 0 & $0.10 \times 0 = 0.00$ \\
EX to 1st \& 2nd   & 0 & $0.10 \times 0 = 0.00$ \\
MEM to 1st only    & 1 & $0.00 \times 1 = 0.00$ \\
MEM to 2nd only    & 0 & $0.05 \times 0 = 0.00$ \\
MEM to 1st \& 2nd  & 1 & $0.05 \times 1 = 0.05$ \\
\midrule
\multicolumn{2}{l}{Total extra CPI} & $0.05$ \\
\bottomrule
\end{tabular}
\end{center}
\[
  \text{CPI}_{\text{full fwd}} = 1 + 0.05 = 1.05
\]

Percentage of cycles stalled:
\[
  \frac{0.05}{1.05} \approx 4.76\%
\]

\phantomsection\addcontentsline{toc}{subsubsection}{4.26.6}
\medskip\noindent\textbf{4.26.6}
Speedup of each forwarding configuration relative to no forwarding
(assuming the same clock cycle time):

\[
  \text{Speedup}_{\text{EX-EX}} = \frac{\text{CPI}_{\text{no fwd}}}
    {\text{CPI}_{\text{EX-EX}}}
  = \frac{1.65}{1.15} \approx 1.43
\]

\[
  \text{Speedup}_{\text{MEM-EX}} = \frac{1.65}{1.25} = 1.32
\]

\[
  \text{Speedup}_{\text{full}} = \frac{1.65}{1.05} \approx 1.57
\]

EX-EX forwarding alone provides a larger benefit than MEM-EX forwarding alone
because EX-to-1st dependencies (ALU $\to$ next instruction) are far more
common than load-use dependencies. Full forwarding provides the greatest
speedup.

\phantomsection\addcontentsline{toc}{subsubsection}{4.26.7}
\medskip\noindent\textbf{4.26.7}
``Time-travel'' forwarding eliminates \emph{all} data hazards, so
$\text{CPI}_{\text{time-travel}} = 1.0$. However, the additional circuitry
adds 100\,ps to the EX stage.

Let the original pipeline stage latencies be as given in Problem~4.16:
IF $= 250$\,ps, ID $= 350$\,ps, EX $= 150$\,ps, MEM $= 300$\,ps,
WB $= 200$\,ps. The clock cycle is determined by the slowest stage:
$\text{Clk}_{\text{full fwd}} = 350$\,ps (ID stage).

With time-travel forwarding, the EX stage becomes $150 + 100 = 250$\,ps. The
bottleneck remains the ID stage at 350\,ps, so the clock cycle is unchanged:
$\text{Clk}_{\text{time-travel}} = 350$\,ps.

\[
  \text{Speedup over full forwarding}
  = \frac{\text{CPI}_{\text{full fwd}} \times \text{Clk}_{\text{full fwd}}}
         {\text{CPI}_{\text{time-travel}} \times \text{Clk}_{\text{time-travel}}}
  = \frac{1.05 \times 350}{1.0 \times 350}
  = 1.05
\]

The additional speedup is $\approx 5\%$ over the fastest previous design.

\phantomsection\addcontentsline{toc}{subsubsection}{4.26.8}
\medskip\noindent\textbf{4.26.8}
The hazard table has separate entries for ``EX to 1st'' and ``EX to 1st and EX
to 2nd,'' but \emph{does} include ``MEM to 1st and MEM to 2nd.'' A question
arises: why is there no cross-type combined entry such as ``EX to 1st and MEM
to 2nd'' or ``MEM to 1st and EX to 2nd''?

The reason is that these cross-type combinations cannot occur from a
\emph{single} producing instruction. A given instruction produces its result
in \emph{either} EX (ALU/R-type/I-type instructions) \emph{or} MEM (load
instructions), but not both. Therefore:

\begin{itemize}
  \item If the same instruction creates both a ``to 1st'' and a ``to 2nd''
    dependency, both dependencies originate from the same pipeline stage
    (either both EX or both MEM).
  \item ``EX to 1st and MEM to 2nd'' would require the same producing
    instruction to generate its result in EX for one consumer and in MEM for
    another, which is impossible.
\end{itemize}

Hence the only meaningful combined entries are ``EX to 1st \& EX to 2nd'' and
``MEM to 1st \& MEM to 2nd.''

%-----------------------------------------------------------------------------
\subsection*{4.27 --- Hazard Detection Code Analysis}
\hypertarget{sol-4.27}{}
\markright{4.27}
\addcontentsline{toc}{subsection}{4.27 --- Hazard Detection Code Analysis}

Reference code sequence on a five-stage pipeline:
\begin{verbatim}
    add $s3, $s1, $s4     # I1
    lw  $s2, 4($s5)       # I2
    or  $s2, $s4, $s1     # I3
    sw  $s2, 0($s5)       # I4
    add $s4, $s2, $s3     # I5
\end{verbatim}

\phantomsection\addcontentsline{toc}{subsubsection}{4.27.1}
\medskip\noindent\textbf{4.27.1}
\textbf{Situation:} The processor has forwarding paths but the hazard detection
unit is missing (not implemented).

First, identify all data dependencies and determine whether any require the
hazard detection unit (i.e., load-use hazards):

\begin{center}
\begin{tabular}{llll}
\toprule
\textbf{Producer} & \textbf{Consumer} & \textbf{Register} & \textbf{Type} \\
\midrule
I1 (\texttt{add \$s3}) & I5 (\texttt{add}) & \texttt{\$s3} &
  EX to 4th (no hazard) \\
I3 (\texttt{or \$s2})  & I4 (\texttt{sw})  & \texttt{\$s2} &
  EX to 1st (forwarding) \\
I3 (\texttt{or \$s2})  & I5 (\texttt{add}) & \texttt{\$s2} &
  EX to 2nd (forwarding) \\
\bottomrule
\end{tabular}
\end{center}

Key observation: I2 (\texttt{lw \$s2}) writes \texttt{\$s2}, but I3
(\texttt{or \$s2, \$s4, \$s1}) also writes \texttt{\$s2} and does \emph{not}
read \texttt{\$s2} --- it reads \texttt{\$s4} and \texttt{\$s1}. Therefore I2's
result for \texttt{\$s2} is never consumed by any later instruction (I3
overwrites it). There is no load-use hazard in this sequence.

\textbf{Conclusion:} Even without the hazard detection unit, this code executes
correctly. All data dependencies are resolved by EX-EX or MEM-EX forwarding
paths. No stalls are needed because there is no load-use hazard (the only
case that requires the hazard detection unit to insert a bubble).

\phantomsection\addcontentsline{toc}{subsubsection}{4.27.2}
\medskip\noindent\textbf{4.27.2}
\textbf{Goal:} Rearrange/modify the code to minimize NOPs, assuming a pipeline
with \emph{no} forwarding and \emph{no} hazard detection. Register
\texttt{\$s6} is available as a temporary.

Without forwarding, a consumer must be at least 2 instructions after the
producer (so the producer completes WB while the consumer is in ID). We
identify dependencies in the original code:

\begin{itemize}
  \item I3 (\texttt{or}) $\to$ I4 (\texttt{sw}): gap = 1, need 2 more $\to$ 2 NOPs
  \item I3 (\texttt{or}) $\to$ I5 (\texttt{add}): gap = 2, need 0 more (if 2 NOPs
    already inserted after I3 for I4, I5 is also satisfied)
\end{itemize}

All other dependencies have sufficient separation. Original code needs 2 NOPs:
\begin{verbatim}
    add $s3, $s1, $s4     # I1
    lw  $s2, 4($s5)       # I2 (result overwritten by I3)
    or  $s2, $s4, $s1     # I3
    nop                    # stall 1 (wait for $s2)
    nop                    # stall 2
    sw  $s2, 0($s5)       # I4: reads $s2 from I3
    add $s4, $s2, $s3     # I5: reads $s2 from I3, $s3 from I1
\end{verbatim}

We can reorder instructions to fill the NOP slots. Since I1 is independent of
I3, and I2's result is overwritten by I3, we can move I1 and I2 after I3:
\begin{verbatim}
    or  $s2, $s4, $s1     # I3 (moved first, no input dependency)
    add $s3, $s1, $s4     # I1 (independent; fills NOP slot)
    lw  $s6, 4($s5)       # I2' (use $s6 since $s2 overwritten anyway)
    sw  $s2, 0($s5)       # I4: reads $s2 from I3 (3 instr gap -- OK)
    add $s4, $s2, $s3     # I5: reads $s2 (4 gap), $s3 (3 gap -- OK)
\end{verbatim}

However, we must verify that the original program semantics are preserved. In
the original code, \texttt{lw \$s2, 4(\$s5)} loads a value that is immediately
overwritten by \texttt{or}. If the load was not used elsewhere (which it is not
in this sequence), we can safely rename it to \texttt{\$s6} or even eliminate
it if it has no side effects needed later. Assuming the load's result may be
needed after this sequence, we rename it to \texttt{\$s6}:

\begin{verbatim}
    or  $s2, $s4, $s1     # produces $s2
    add $s3, $s1, $s4     # produces $s3 (no dep on $s2)
    lw  $s6, 4($s5)       # loads to $s6 (renamed from $s2)
    sw  $s2, 0($s5)       # reads $s2 (3-instr gap, OK)
    add $s4, $s2, $s3     # reads $s2 and $s3 (gaps >= 2, OK)
\end{verbatim}

\textbf{Result:} 0 NOPs needed after reordering. $\checkmark$

\phantomsection\addcontentsline{toc}{subsubsection}{4.27.3}
\medskip\noindent\textbf{4.27.3}
Pipeline execution with forwarding for the first seven cycles. We show which
forwarding and hazard detection signals are asserted.

Pipeline diagram (with forwarding, no stalls needed):
\begin{verbatim}
              CC1  CC2  CC3  CC4  CC5  CC6  CC7  CC8  CC9
  I1: add    IF   ID   EX   MEM  WB
  I2: lw          IF   ID   EX   MEM  WB
  I3: or               IF   ID   EX   MEM  WB
  I4: sw                    IF   ID   EX   MEM  WB
  I5: add                        IF   ID   EX   MEM  WB
\end{verbatim}

Forwarding signals asserted each cycle:

\begin{center}
\begin{tabular}{cl}
\toprule
\textbf{Cycle} & \textbf{Signals Asserted} \\
\midrule
CC1 & None (I1 in IF) \\
CC2 & None (I1 in ID, I2 in IF) \\
CC3 & None (I1 in EX, I2 in ID, I3 in IF --- no dependencies yet) \\
CC4 & None (I1 in MEM, I2 in EX, I3 in ID --- I3 reads \$s4, \$s1; \\
    & \quad no producer in EX/MEM or MEM/WB writes those registers) \\
CC5 & \texttt{ForwardA = 00, ForwardB = 00} for I4 (sw reads \$s2 from I3; \\
    & \quad I3 is in MEM/WB $\to$ \textbf{MEM-EX forward} for \$s2) \\
    & \quad Specifically: MEM/WB.Rd = \$s2 matches ID/EX.Rs of I4 \\
CC6 & I5 in ID reads \$s2 and \$s3; I3 in WB (done), I4 in MEM/WB: \\
    & \quad I3.Rd = \$s2 forwarded from MEM/WB to I5? No, I4 is \texttt{sw}, \\
    & \quad no write. Check EX/MEM: I4 is sw (RegWrite=0). \\
    & \quad \$s2: I3 already completed WB, value in register file. \\
    & \quad \$s3: I1 completed WB in CC5, value in register file. \\
    & \quad No forwarding needed (values available from reg file). \\
CC7 & I5 in EX. No active forwarding. \\
\bottomrule
\end{tabular}
\end{center}

Hazard detection unit: No load-use hazard detected in any cycle (I2 produces
\texttt{\$s2} but I3 does not read \texttt{\$s2}), so the hazard detection unit
never asserts a stall signal.

Summary of forwarding during CC5: \texttt{ForwardA} or \texttt{ForwardB} is set
to select the MEM/WB forwarding path to deliver \texttt{\$s2} from I3's result
to I4's \texttt{sw} data input.

\phantomsection\addcontentsline{toc}{subsubsection}{4.27.4}
\medskip\noindent\textbf{4.27.4}
Without forwarding, the hazard detection unit must detect \emph{all} RAW
hazards (not just load-use). It needs to compare the destination register of
instructions in EX, MEM, and WB stages against the source registers of the
instruction in ID.

\textbf{New inputs needed:}
\begin{itemize}
  \item \texttt{EX/MEM.RegisterRd} and \texttt{EX/MEM.RegWrite} --- to detect
    dependencies from the instruction 2 stages ahead.
  \item \texttt{MEM/WB.RegisterRd} and \texttt{MEM/WB.RegWrite} --- to detect
    dependencies from the instruction 3 stages ahead (already partially handled
    by write-first/read-second, but without forwarding we must stall if the
    write hasn't completed).
\end{itemize}

In addition to the existing inputs (\texttt{ID/EX.RegisterRt} for load-use
detection), we need:
\begin{itemize}
  \item \texttt{IF/ID.RegisterRs} and \texttt{IF/ID.RegisterRt} --- source
    registers of the instruction currently in ID.
  \item \texttt{ID/EX.RegisterRd} and \texttt{ID/EX.RegWrite} --- destination
    of instruction in EX.
\end{itemize}

\textbf{New output:} The stall/bubble signal (same as before), but it must now
be asserted for any RAW hazard, not just load-use.

\textbf{Example from the code:} I3 (\texttt{or}) writes \texttt{\$s2}. I4
(\texttt{sw}) reads \texttt{\$s2} in the very next cycle. Without forwarding,
the hazard detection unit must stall I4 for 2 cycles until I3's result is
written back.

\phantomsection\addcontentsline{toc}{subsubsection}{4.27.5}
\medskip\noindent\textbf{4.27.5}
Without forwarding, we trace the first 5 cycles and identify stalls. The
hazard detection unit inserts bubbles whenever a source register in ID matches
a destination register in EX or MEM (and that instruction writes a register).

\begin{verbatim}
         CC1  CC2  CC3  CC4  CC5  CC6  CC7  CC8  CC9
  I1:    IF   ID   EX   MEM  WB
  I2:         IF   ID   EX   MEM  WB
  I3:              IF   ID   EX   MEM  WB
  I4:                   IF   ID*  ID*  ID   EX   MEM  WB
  I5:                        IF   IF   IF   ID   ...
\end{verbatim}

During the first 5 cycles (CC1--CC5):

\begin{itemize}
  \item CC1--CC3: No stalls. I1, I2, I3 flow through IF, ID, EX without
    dependencies among them (I3 reads \texttt{\$s4} and \texttt{\$s1}; I1 writes
    \texttt{\$s3}, I2 writes \texttt{\$s2} --- no match with I3's sources).
  \item CC4: I4 enters ID. I4 (\texttt{sw \$s2}) reads \texttt{\$s2}. I3 is
    in EX and writes \texttt{\$s2}. Hazard detected $\to$ \textbf{stall}.
  \item CC5: I4 still in ID. I3 now in MEM, still writing \texttt{\$s2} (not
    yet in WB). Hazard persists $\to$ \textbf{stall}.
\end{itemize}

\textbf{Result:} Approximately 2 stall cycles occur within the first 5
pipeline cycles (at CC4 and CC5), both due to the I3 $\to$ I4 dependency on
\texttt{\$s2}.

\phantomsection\addcontentsline{toc}{subsubsection}{4.27.6}
\medskip\noindent\textbf{4.27.6}
Total stall cycles for the entire code sequence without forwarding:

Dependencies requiring stalls:
\begin{enumerate}
  \item I3 (\texttt{or \$s2}) $\to$ I4 (\texttt{sw \$s2}): gap = 1 instruction,
    need gap of 2 $\to$ \textbf{2 stall cycles}.
  \item I3 (\texttt{or \$s2}) $\to$ I5 (\texttt{add \$s4, \$s2, \$s3}): original
    gap = 2 instructions, but after inserting 2 stalls for I4, the gap between
    I3 and I5 becomes 4 $\to$ already satisfied, \textbf{0 additional stalls}.
  \item I1 (\texttt{add \$s3}) $\to$ I5 (\texttt{add \$s4, \$s2, \$s3}): original
    gap = 4 instructions $\to$ \textbf{0 stalls}.
\end{enumerate}

\textbf{Total stall cycles for the entire sequence: 2.} $\checkmark$

%-----------------------------------------------------------------------------
\subsection*{4.28 --- Branch Prediction CPI Impact}
\hypertarget{sol-4.28}{}
\markright{4.28}
\addcontentsline{toc}{subsection}{4.28 --- Branch Prediction CPI Impact}

Instruction mix and predictor accuracy:

\begin{center}
\begin{tabular}{ccccc}
\toprule
\textbf{R-type} & \textbf{beq/bne} & \textbf{jal} & \textbf{lw} &
\textbf{sw} \\
\midrule
40\% & 25\% & 5\% & 20\% & 10\% \\
\bottomrule
\end{tabular}
\end{center}

\begin{center}
\begin{tabular}{ccc}
\toprule
\textbf{Always-Taken} & \textbf{Always-Not-Taken} & \textbf{2-Bit} \\
\midrule
45\% & 55\% & 85\% \\
\bottomrule
\end{tabular}
\end{center}

Assumptions: branches (\texttt{beq}/\texttt{bne}) are resolved in the ID
stage (with dedicated comparison hardware, as in Section~4.9), so the
misprediction penalty is \textbf{1 cycle} (one instruction fetched after the
branch must be flushed).

\phantomsection\addcontentsline{toc}{subsubsection}{4.28.1}
\medskip\noindent\textbf{4.28.1}
Extra CPI due to mispredicted branches with the \textbf{always-taken}
predictor:
\[
  \text{Extra CPI} = \text{branch fraction} \times \text{misprediction rate}
    \times \text{penalty}
\]
\[
  = 0.25 \times (1 - 0.45) \times 1
  = 0.25 \times 0.55 \times 1
  = \boxed{0.1375}
\]

\phantomsection\addcontentsline{toc}{subsubsection}{4.28.2}
\medskip\noindent\textbf{4.28.2}
Extra CPI with the \textbf{always-not-taken} predictor:
\[
  \text{Extra CPI} = 0.25 \times (1 - 0.55) \times 1
  = 0.25 \times 0.45 \times 1
  = \boxed{0.1125}
\]

\phantomsection\addcontentsline{toc}{subsubsection}{4.28.3}
\medskip\noindent\textbf{4.28.3}
Extra CPI with the \textbf{2-bit} predictor:
\[
  \text{Extra CPI} = 0.25 \times (1 - 0.85) \times 1
  = 0.25 \times 0.15 \times 1
  = \boxed{0.0375}
\]

\phantomsection\addcontentsline{toc}{subsubsection}{4.28.4}
\medskip\noindent\textbf{4.28.4}
With the 2-bit predictor, the CPI is:
\[
  \text{CPI}_{\text{original}} = 1 + 0.0375 = 1.0375
\]

If we could convert \emph{all} branch instructions to ALU instructions
(eliminating all mispredictions), the branch fraction becomes 0\%, so:
\[
  \text{CPI}_{\text{no branches}} = 1.0
\]

Speedup:
\[
  \text{Speedup} = \frac{1.0375}{1.0} = \boxed{1.0375}
\]

This is a modest $\approx 3.75\%$ improvement, because the 2-bit predictor
is already quite accurate.

\phantomsection\addcontentsline{toc}{subsubsection}{4.28.5}
\medskip\noindent\textbf{4.28.5}
If we convert \emph{half} of the branch instructions to (non-ALU)
instructions that do not cause mispredictions, the branch fraction becomes:
\[
  \text{New branch fraction} = \frac{25\%}{2} = 12.5\%
\]

New extra CPI:
\[
  \text{Extra CPI}_{\text{new}} = 0.125 \times 0.15 \times 1 = 0.01875
\]
\[
  \text{CPI}_{\text{new}} = 1 + 0.01875 = 1.01875
\]

Speedup over the original (with 2-bit predictor):
\[
  \text{Speedup} = \frac{1.0375}{1.01875} \approx \boxed{1.018}
\]

A modest $\approx 1.8\%$ improvement.

\phantomsection\addcontentsline{toc}{subsubsection}{4.28.6}
\medskip\noindent\textbf{4.28.6}
Given that 80\% of all branch instructions are ``easy-to-predict'' loop
branches that are always predicted correctly, and the overall 2-bit predictor
accuracy is 85\%, we find the accuracy on the remaining 20\% of branches.

Let $x$ be the accuracy on the hard-to-predict 20\%:
\[
  0.80 \times 1.00 + 0.20 \times x = 0.85
\]
\[
  0.80 + 0.20x = 0.85
\]
\[
  0.20x = 0.05
\]
\[
  x = \frac{0.05}{0.20} = \boxed{0.25 = 25\%}
\]

The 2-bit predictor achieves only 25\% accuracy on the remaining 20\% of
(hard-to-predict) branch instructions. This is worse than random guessing
(50\%), indicating that these branches have patterns that actively confuse the
2-bit predictor.

%-----------------------------------------------------------------------------
\subsection*{4.29 --- Branch Predictor Design}
\hypertarget{sol-4.29}{}
\markright{4.29}
\addcontentsline{toc}{subsection}{4.29 --- Branch Predictor Design}

Repeating branch pattern: T, T, NT, T, T, NT, \ldots\ (period~3).

\phantomsection\addcontentsline{toc}{subsubsection}{4.29.1}
\medskip\noindent\textbf{4.29.1}
\textbf{Always-Taken predictor:} predicts T every time.
\begin{itemize}
  \item Correct on T outcomes: 2 out of every 3 branches.
  \item Incorrect on NT outcomes: 1 out of every 3 branches.
\end{itemize}
\[
  \text{Accuracy}_{\text{always-T}} = \frac{2}{3} \approx \boxed{66.7\%}
\]

\textbf{Always-Not-Taken predictor:} predicts NT every time.
\begin{itemize}
  \item Correct on NT outcomes: 1 out of every 3 branches.
  \item Incorrect on T outcomes: 2 out of every 3 branches.
\end{itemize}
\[
  \text{Accuracy}_{\text{always-NT}} = \frac{1}{3} \approx \boxed{33.3\%}
\]

\phantomsection\addcontentsline{toc}{subsubsection}{4.29.2}
\medskip\noindent\textbf{4.29.2}
2-bit predictor states (from Figure~4.63):
\begin{itemize}
  \item State 0: Strongly Not Taken (predict NT)
  \item State 1: Weakly Not Taken (predict NT)
  \item State 2: Weakly Taken (predict T)
  \item State 3: Strongly Taken (predict T)
\end{itemize}

Transitions: actual outcome T increments the state (up to 3); actual outcome
NT decrements (down to 0). Starting state: 0 (Strongly Not Taken).

\begin{center}
\begin{tabular}{ccccc}
\toprule
\textbf{Branch \#} & \textbf{State before} & \textbf{Prediction} &
  \textbf{Actual} & \textbf{Correct?} \\
\midrule
1 & 0 (Strongly NT) & NT & T  & No  \\
2 & 1 (Weakly NT)   & NT & T  & No  \\
3 & 2 (Weakly T)    & T  & NT & No  \\
4 & 1 (Weakly NT)   & NT & T  & No  \\
\bottomrule
\end{tabular}
\end{center}

After the state transitions: $0 \xrightarrow{T} 1 \xrightarrow{T} 2
\xrightarrow{NT} 1 \xrightarrow{T} 2$.

\[
  \text{Accuracy (first 4 branches)} = \frac{0}{4} = \boxed{0\%}
\]

\phantomsection\addcontentsline{toc}{subsubsection}{4.29.3}
\medskip\noindent\textbf{4.29.3}
Continuing the trace from where 4.29.2 left off (state = 2 after branch 4):

\begin{center}
\begin{tabular}{ccccc}
\toprule
\textbf{Branch \#} & \textbf{State before} & \textbf{Prediction} &
  \textbf{Actual} & \textbf{Correct?} \\
\midrule
5  & 2 (Weakly T)    & T & T  & Yes \\
6  & 3 (Strongly T)  & T & NT & No  \\
7  & 2 (Weakly T)    & T & T  & Yes \\
8  & 3 (Strongly T)  & T & T  & Yes \\
9  & 3 (Strongly T)  & T & NT & No  \\
10 & 2 (Weakly T)    & T & T  & Yes \\
11 & 3 (Strongly T)  & T & T  & Yes \\
12 & 3 (Strongly T)  & T & NT & No  \\
\bottomrule
\end{tabular}
\end{center}

Starting from branch~7 onward, the pattern of states and outcomes repeats
with a period of 3:
\begin{itemize}
  \item State 2 $\xrightarrow{T}$ 3 (correct), State 3 $\xrightarrow{T}$ 3
    (correct), State 3 $\xrightarrow{NT}$ 2 (incorrect).
\end{itemize}

In steady state, the predictor always predicts T (from states 2 or 3) and is
correct on 2 out of every 3 branches (the two T outcomes) and incorrect on
the 1 NT outcome.

\[
  \text{Steady-state accuracy} = \frac{2}{3} \approx \boxed{66.7\%}
\]

This matches the always-taken predictor's accuracy. The 2-bit predictor cannot
do better than always-taken for this pattern because after warming up, it
effectively becomes an always-taken predictor (stuck in states 2--3).

\phantomsection\addcontentsline{toc}{subsubsection}{4.29.4}
\medskip\noindent\textbf{4.29.4}
\textbf{Design a predictor with perfect accuracy for the repeating pattern
T, T, NT.}

We use a \textbf{3-state counter} (mod~3) that cycles unconditionally. The
predictor has no input other than the clock and a signal indicating a
conditional branch is being executed.

\textbf{States and outputs:}
\begin{center}
\begin{tabular}{ccc}
\toprule
\textbf{State} & \textbf{Prediction} & \textbf{Next State} \\
\midrule
$S_0$ & 1 (Taken)     & $S_1$ \\
$S_1$ & 1 (Taken)     & $S_2$ \\
$S_2$ & 0 (Not Taken) & $S_0$ \\
\bottomrule
\end{tabular}
\end{center}

\textbf{Implementation:} This can be built with a 2-bit register (encoding
3 states) and simple combinational logic:
\begin{itemize}
  \item Output = $\overline{S_1 \cdot S_0}$ (output is 0 only when state =
    $S_2 = \texttt{10}_2$, otherwise 1).
  \item Or more precisely, if we encode $S_0 = 00$, $S_1 = 01$, $S_2 = 10$:
    \begin{itemize}
      \item Prediction output = NOT(bit$_1$) (the high bit)
      \item Next state: $00 \to 01 \to 10 \to 00$ (a standard mod-3 counter)
    \end{itemize}
  \item The counter advances by 1 (mod~3) on every branch execution.
\end{itemize}

Starting in $S_0$, the predictions are: T, T, NT, T, T, NT, \ldots\ which
exactly matches the pattern. \textbf{Accuracy = 100\%.} $\checkmark$

\phantomsection\addcontentsline{toc}{subsubsection}{4.29.5}
\medskip\noindent\textbf{4.29.5}
The \textbf{opposite pattern} is: NT, NT, T, NT, NT, T, \ldots

Our mod-3 predictor (from 4.29.4) still outputs: T, T, NT, T, T, NT, \ldots

Comparing prediction vs.\ actual outcome:

\begin{center}
\begin{tabular}{cccc}
\toprule
\textbf{Branch \#} & \textbf{Prediction} & \textbf{Actual} &
  \textbf{Correct?} \\
\midrule
1 & T  & NT & No \\
2 & T  & NT & No \\
3 & NT & T  & No \\
4 & T  & NT & No \\
$\vdots$ & $\vdots$ & $\vdots$ & $\vdots$ \\
\bottomrule
\end{tabular}
\end{center}

Every prediction is the opposite of the actual outcome.

\[
  \text{Accuracy} = \boxed{0\%}
\]

The predictor from 4.29.4 achieves the worst possible accuracy on the opposite
pattern, because the two patterns are exact complements of each other.

\phantomsection\addcontentsline{toc}{subsubsection}{4.29.6}
\medskip\noindent\textbf{4.29.6}
\textbf{Design a predictor that eventually achieves perfect accuracy on both
the pattern T, T, NT and its opposite NT, NT, T.}

The predictor now receives an additional input: the \textbf{actual branch
outcome} after each branch resolves. This feedback allows it to determine which
of the two patterns it is observing.

\textbf{Design: mod-3 counter with a mode bit.}

\textbf{State:} Two components:
\begin{enumerate}
  \item A 2-bit \textbf{phase counter} $P \in \{0, 1, 2\}$ cycling mod~3
    (same as 4.29.4).
  \item A 1-bit \textbf{mode register} $M \in \{0, 1\}$:
    \begin{itemize}
      \item $M = 0$ (normal): predictions are T, T, NT (for phases 0, 1, 2).
      \item $M = 1$ (inverted): predictions are NT, NT, T (for phases 0, 1, 2).
    \end{itemize}
\end{enumerate}

\textbf{Operation on each branch:}
\begin{enumerate}
  \item Read the current phase $P$ and mode $M$.
  \item Compute prediction:
    \[
      \text{Pred} =
      \begin{cases}
        \text{base}(P) & \text{if } M = 0 \\
        \overline{\text{base}(P)} & \text{if } M = 1
      \end{cases}
    \]
    where $\text{base}(0) = \text{base}(1) = \text{T}$ and
    $\text{base}(2) = \text{NT}$.
  \item After the actual outcome is known, compare prediction with actual:
    \begin{itemize}
      \item If \textbf{mismatch}: toggle $M$ (flip the mode bit).
      \item If \textbf{match}: keep $M$ unchanged.
    \end{itemize}
  \item Advance phase: $P \leftarrow (P + 1) \bmod 3$.
\end{enumerate}

\textbf{Total state:} 2 bits for the phase counter + 1 bit for the mode
= \textbf{3 bits}.

\textbf{Behavior on pattern T, T, NT} (starting with $M = 0$, $P = 0$):

\begin{center}
\begin{tabular}{cccccc}
\toprule
\textbf{\#} & $P$ & $M$ & \textbf{Pred} & \textbf{Actual} &
  \textbf{Action} \\
\midrule
1 & 0 & 0 & T  & T  & match; $M$ stays 0 \\
2 & 1 & 0 & T  & T  & match; $M$ stays 0 \\
3 & 2 & 0 & NT & NT & match; $M$ stays 0 \\
\bottomrule
\end{tabular}
\end{center}

Perfect from the start. $\checkmark$

\textbf{Behavior on opposite pattern NT, NT, T} (starting with $M = 0$, $P = 0$):

\begin{center}
\begin{tabular}{cccccc}
\toprule
\textbf{\#} & $P$ & $M$ & \textbf{Pred} & \textbf{Actual} &
  \textbf{Action} \\
\midrule
1 & 0 & 0 & T  & NT & mismatch; toggle $M \to 1$ \\
2 & 1 & 1 & NT & NT & match; $M$ stays 1 \\
3 & 2 & 1 & T  & T  & match; $M$ stays 1 \\
4 & 0 & 1 & NT & NT & match; $M$ stays 1 \\
5 & 1 & 1 & NT & NT & match; $M$ stays 1 \\
6 & 2 & 1 & T  & T  & match; $M$ stays 1 \\
\bottomrule
\end{tabular}
\end{center}

After only 1 misprediction (branch~1), the mode toggles and the predictor
locks into the correct mode. From branch~2 onward, all predictions are
correct.

\textbf{Steady-state accuracy on both patterns: 100\%.} $\checkmark$

\textbf{Hardware summary:}
\begin{itemize}
  \item 2-bit mod-3 counter for phase ($P$).
  \item 1-bit mode register ($M$).
  \item XOR gate to invert the base prediction when $M = 1$.
  \item XOR comparator between prediction and actual outcome to determine
    whether to toggle $M$.
  \item Mod-3 increment logic (same as 4.29.4).
\end{itemize}
%-----------------------------------------------------------------------------
\subsection*{4.30 --- Exception Handling in Pipeline}
\hypertarget{sol-4.30}{}
\markright{4.30}
\addcontentsline{toc}{subsection}{4.30 --- Exception Handling in Pipeline}

The following two instructions are used for the first three sub-parts:

\begin{center}
\begin{tabular}{cc}
\textbf{Instruction 1} & \textbf{Instruction 2} \\
\texttt{beq \$s1, \$s0, label} & \texttt{add \$s3, \$s1, \$s2}
\end{tabular}
\end{center}

\phantomsection\addcontentsline{toc}{subsubsection}{4.30.1}
\medskip\noindent\textbf{4.30.1}
Each instruction can trigger exceptions at particular pipeline stages:

\begin{center}
\begin{tabular}{lll}
\toprule
\textbf{Instruction} & \textbf{Exception} & \textbf{Stage Detected} \\
\midrule
\texttt{beq} & Invalid/undefined instruction (bad opcode) & ID \\
\texttt{beq} & Instruction address exception (PC misaligned) & IF \\
\midrule
\texttt{add} & Invalid/undefined instruction (bad opcode or funct) & ID \\
\texttt{add} & Arithmetic overflow & EX \\
\texttt{add} & Instruction address exception (PC misaligned) & IF \\
\bottomrule
\end{tabular}
\end{center}

\noindent Key distinction: \texttt{beq} performs a subtraction in the ALU to compare
registers, but MIPS does not generate an arithmetic overflow exception for
branch comparisons.  The \texttt{add} instruction, however, \emph{does}
generate an overflow exception when the signed result cannot be represented in
32 bits.

\phantomsection\addcontentsline{toc}{subsubsection}{4.30.2}
\medskip\noindent\textbf{4.30.2}
To support a separate handler address for each exception type, the pipeline
organisation must be modified as follows:

\begin{enumerate}
\item \textbf{Exception cause encoder.}  Add logic that identifies the
      exception type and produces a cause code (e.g., 2-bit encoding:
      00 = undefined instruction, 01 = overflow, 10 = address exception).

\item \textbf{Handler-address selection.}  Replace the single hardwired
      exception handler address with a \emph{mux} (or lookup table) that maps
      the cause code to the corresponding handler address.  The mux output
      feeds the PC input.

\item \textbf{EPC register.}  Add the Exception Program Counter register.
      When an exception is detected, the PC of the faulting instruction is
      saved in EPC so the handler can return to the correct location.

\item \textbf{Pipeline flush logic.}  When an exception is detected in any
      stage, all instructions that entered the pipeline \emph{after} the
      faulting instruction must be flushed (their control signals are zeroed).

\item \textbf{Priority logic.}  If multiple exceptions occur simultaneously
      (e.g., in different pipeline stages), the exception from the earliest
      instruction (furthest along in the pipeline) is handled first.
\end{enumerate}

\noindent The modified next-PC selection becomes:
\[
\text{Next PC} = \begin{cases}
  \text{PC} + 4 & \text{normal} \\
  \text{branch target} & \text{branch taken} \\
  \text{Handler}_i & \text{exception of type } i
\end{cases}
\]

\phantomsection\addcontentsline{toc}{subsubsection}{4.30.3}
\medskip\noindent\textbf{4.30.3}
Suppose \texttt{beq} (I1) causes an exception---for example, an undefined
instruction detected in the ID stage.  At that point the second instruction
\texttt{add} (I2) has already been fetched and is in the IF stage.  The
pipeline must:

\begin{enumerate}
\item Detect the exception at the end of the ID stage for I1.
\item Save the PC of I1 in the EPC register.
\item Flush I2 by converting it to a \texttt{nop} (zero all its control
      signals).
\item Flush I1's remaining pipeline stages (prevent I1 from writing results).
\item Load the appropriate exception handler address into the PC.
\end{enumerate}

\noindent Pipeline execution diagram (exception detected in ID of \texttt{beq}):

\begin{verbatim}
             CC1    CC2    CC3    CC4    CC5    CC6    CC7
    beq:     IF     ID*    ---    ---    ---
    add:            IF     ---    ---    ---
    handler:               IF     ID     EX     MEM    WB
\end{verbatim}

\noindent \texttt{*} = exception detected. ``---'' indicates a flushed stage.
The handler's first instruction begins fetching in CC3.

If instead the exception occurs later (e.g., \texttt{add} overflows in EX):

\begin{verbatim}
             CC1    CC2    CC3    CC4    CC5    CC6    CC7    CC8
    beq:     IF     ID     EX     MEM    WB
    add:            IF     ID     EX*    ---    ---
    (next):                IF     ID     ---    ---
    handler:                             IF     ID     EX     MEM
\end{verbatim}

\noindent Here \texttt{add}'s overflow is detected at the end of EX (CC4).
Any instruction fetched after \texttt{add} is flushed, and the handler begins
fetching in CC5.

\phantomsection\addcontentsline{toc}{subsubsection}{4.30.4}
\medskip\noindent\textbf{4.30.4}
In vectored exception handling with the vector table stored in data memory:

\begin{enumerate}
\item The exception type generates an \emph{index} (cause code).
\item A base address for the vector table is hardwired (e.g., \texttt{0x80000000}).
\item The handler address is read from memory at
      $\text{base} + \text{cause code} \times 4$.
\item This requires an \emph{additional memory access} to read the handler
      address from data memory before the handler can begin executing.
\end{enumerate}

\noindent Required hardware modifications:
\begin{itemize}
\item An adder to compute $\text{table base} + \text{cause} \times 4$.
\item A path from data memory output back to the PC input mux (or a dedicated
      read port).
\item The pipeline must stall for one extra cycle while the vector table entry
      is fetched from memory.
\end{itemize}

\noindent Pipeline diagram (exception in ID of \texttt{beq}, vectored via memory):

\begin{verbatim}
             CC1    CC2    CC3    CC4    CC5    CC6    CC7    CC8
    beq:     IF     ID*    ---    ---    ---    ---
    add:            IF     ---    ---    ---    ---
    (vec rd):              MEM    ---                         [read table]
    handler:                      IF     ID     EX     MEM    WB
\end{verbatim}

\noindent CC3 is used to read the handler address from the vector table in
data memory.  The actual handler begins fetching in CC4, one cycle later
than the non-vectored case.

\phantomsection\addcontentsline{toc}{subsubsection}{4.30.5}
\medskip\noindent\textbf{4.30.5}
To emulate vectored exception handling with only a single fixed handler
address, the code at the fixed address must read the Cause register,
determine the exception type, and branch to the appropriate specific handler.
A jump-table approach is efficient:

\begin{verbatim}
    # Code at the single fixed exception handler address
    mfc0  $k0, $13            # Read Cause register
    andi  $k0, $k0, 0x7C     # Extract ExcCode field (bits 6:2)
                              #   already a word offset (x4)
    la    $k1, jump_table     # Load base of handler address table
    add   $k1, $k1, $k0      # Index into table
    lw    $k1, 0($k1)        # Load specific handler address
    jr    $k1                 # Jump to the specific handler

    # Handler address table (in data segment)
    jump_table:
        .word handler_int     # ExcCode 0: interrupt
        .word handler_tlb     # ExcCode 1: TLB modification
        .word handler_tlbl    # ExcCode 2: TLB load
        .word handler_tlbs    # ExcCode 3: TLB store
        .word handler_addrl   # ExcCode 4: address error (load)
        .word handler_addrs   # ExcCode 5: address error (store)
        ...
        .word handler_ovf     # ExcCode 12: overflow
\end{verbatim}

\noindent The key insight is that the \texttt{ExcCode} field (bits [6:2] of
the Cause register) is already shifted left by 2, so it directly serves as a
byte offset into a word-aligned table.  This approach uses 6 instructions
of overhead before reaching the actual handler, compared to zero overhead
with true hardware-vectored exceptions.

%-----------------------------------------------------------------------------
\subsection*{4.31 --- Two-Issue Processor}
\hypertarget{sol-4.31}{}
\markright{4.31}
\addcontentsline{toc}{subsection}{4.31 --- Two-Issue Processor}

The loop under consideration (steady-state body in bold):

\begin{verbatim}
          li   $s1, 0
          j    ENT
    TOP:  add  $s1, $s2, $s0     # ALU
          lw   $s2, 0($s3)       # Memory
          sub  $s3, $s3, $s4     # ALU
    ENT:  bne  $s0, $s1, TOP     # Branch (ALU slot)
\end{verbatim}

\noindent Two-issue rules: each packet contains at most one memory operation
and one ALU/branch operation.  Full forwarding, perfect branch prediction.
Two instructions in the same packet may not have a dependency between them.

\phantomsection\addcontentsline{toc}{subsubsection}{4.31.1}
\medskip\noindent\textbf{4.31.1}
Analyse the loop body dependencies (steady-state iteration):

\begin{itemize}
\item \texttt{add \$s1, \$s2, \$s0}: reads \texttt{\$s2} (produced by
      \emph{previous} iteration's \texttt{lw}).
\item \texttt{lw \$s2, 0(\$s3)}: reads \texttt{\$s3}, writes \texttt{\$s2}.
\item \texttt{sub \$s3, \$s3, \$s4}: reads \texttt{\$s3} (no dependency on
      \texttt{lw}'s output \texttt{\$s2}).
\item \texttt{bne \$s0, \$s1}: reads \texttt{\$s1} (from \texttt{add}).
\end{itemize}

\noindent Can \texttt{lw} and \texttt{add} be paired?  \texttt{add} reads
\texttt{\$s2} (old value from \emph{previous} \texttt{lw}), and this
\texttt{lw} writes \texttt{\$s2} (new value).  Since register reads happen
before writes in the same cycle, \texttt{add} reads the old \texttt{\$s2}.
The two instructions in this packet have no true (same-iteration) data
dependence, so they \emph{can} issue together.

Can \texttt{sub} and \texttt{bne} be paired?  Both need the ALU/branch slot
--- only one ALU/branch per packet.  They cannot be paired.

\begin{center}
\begin{tabular}{clll}
\toprule
\textbf{Cycle} & \textbf{Memory slot} & \textbf{ALU/Branch slot} & \textbf{Notes} \\
\midrule
1 & \texttt{lw \$s2, 0(\$s3)} & \texttt{add \$s1, \$s2, \$s0} & paired OK \\
2 & \texttt{nop} & \texttt{sub \$s3, \$s3, \$s4} & only ALU \\
3 & \texttt{nop} & \texttt{bne \$s0, \$s1, TOP} & branch; \texttt{\$s1} forwarded from cycle 1 \\
\bottomrule
\end{tabular}
\end{center}

\noindent Pipeline diagram for two iterations (each packet uses both issue
slots; IF/ID/EX/MEM/WB shown for the first instruction in each packet):

\begin{verbatim}
         CC1  CC2  CC3  CC4  CC5  CC6  CC7  CC8  CC9  CC10
Pkt 1a:  IF   ID   EX   MEM  WB                 (lw + add)
Pkt 2a:       IF   ID   EX   MEM  WB            (nop + sub)
Pkt 3a:            IF   ID   EX   MEM  WB       (nop + bne)
Pkt 1b:                 IF   ID   EX   MEM  WB  (lw + add) iter 2
Pkt 2b:                      IF   ID   EX   MEM  WB
Pkt 3b:                           IF   ID   EX   MEM  WB
\end{verbatim}

\noindent Each iteration takes \textbf{3 cycles} for 4 instructions.

\phantomsection\addcontentsline{toc}{subsubsection}{4.31.2}
\medskip\noindent\textbf{4.31.2}
\textbf{One-issue pipeline:} The loop body has 4 instructions.  Check for
data hazards with full forwarding:

\begin{itemize}
\item \texttt{add} $\to$ \texttt{lw}: no dependency (different registers written/read).
\item \texttt{lw} $\to$ \texttt{sub}: \texttt{sub} reads \texttt{\$s3},
      \texttt{lw} reads \texttt{\$s3} --- no RAW hazard (both read).
\item \texttt{sub} $\to$ \texttt{bne}: \texttt{bne} reads \texttt{\$s0, \$s1};
      \texttt{sub} writes \texttt{\$s3} --- no hazard.
\item \texttt{add} $\to$ \texttt{bne}: \texttt{add} writes \texttt{\$s1},
      \texttt{bne} reads \texttt{\$s1}.  \texttt{add} is 3 instructions before
      \texttt{bne}, so forwarding handles it (EX/MEM $\to$ EX, no stall).
\item Cross-iteration: previous \texttt{lw} writes \texttt{\$s2},
      next \texttt{add} reads \texttt{\$s2}.  Separation is 3 instructions
      (\texttt{sub}, \texttt{bne}, then \texttt{add}), so forwarding handles
      it with no stall.
\end{itemize}

One-issue: \textbf{4 cycles} per iteration (no stalls).

\[
\text{Speedup} = \frac{4 \text{ cycles (1-issue)}}{3 \text{ cycles (2-issue)}}
              = \frac{4}{3} \approx 1.33
\]

\phantomsection\addcontentsline{toc}{subsubsection}{4.31.3}
\medskip\noindent\textbf{4.31.3}
To improve performance on the two-issue processor, we want to maximise the
number of instructions that can be paired.  The bottleneck is that
\texttt{sub} and \texttt{bne} both require the ALU slot and cannot be paired.

Reorder the loop body to place \texttt{sub} where it can pair with
\texttt{lw}:

\begin{verbatim}
    TOP:  lw   $s2, 0($s3)       # Memory
          sub  $s3, $s3, $s4     # ALU (no dep on this lw's result)
          add  $s1, $s2, $s0     # ALU (uses $s2 from prev iter's lw)
    ENT:  bne  $s0, $s1, TOP     # Branch
\end{verbatim}

\noindent Correctness check: \texttt{sub} reads \texttt{\$s3} and \texttt{lw}
also reads \texttt{\$s3} (no WAR hazard in hardware); \texttt{sub} does not
use \texttt{\$s2}, so it is independent of \texttt{lw}.  \texttt{add} uses
\texttt{\$s2} from the previous iteration's \texttt{lw}, not this one's, so
the reordering is safe.

\phantomsection\addcontentsline{toc}{subsubsection}{4.31.4}
\medskip\noindent\textbf{4.31.4}
Issue schedule for the optimised code:

\begin{center}
\begin{tabular}{clll}
\toprule
\textbf{Cycle} & \textbf{Memory slot} & \textbf{ALU/Branch slot} & \textbf{Notes} \\
\midrule
1 & \texttt{lw \$s2, 0(\$s3)} & \texttt{sub \$s3, \$s3, \$s4} & no inter-dep \\
2 & \texttt{nop} & \texttt{add \$s1, \$s2, \$s0} & \texttt{\$s2} from prev iter fwd \\
3 & \texttt{nop} & \texttt{bne \$s0, \$s1, TOP} & \texttt{\$s1} fwd from cycle 2 \\
\bottomrule
\end{tabular}
\end{center}

\noindent Still 3 cycles per iteration.  The reordering allowed \texttt{sub}
to pair with \texttt{lw}, but we still have two unpaired ALU operations
(\texttt{add}, \texttt{bne}).  Without additional memory instructions to fill
the memory slot, we cannot do better than 3 cycles without loop unrolling.

\phantomsection\addcontentsline{toc}{subsubsection}{4.31.5}
\medskip\noindent\textbf{4.31.5}
(Same as 4.31.4 --- the pipeline diagram for the optimised code was shown
above.)

\phantomsection\addcontentsline{toc}{subsubsection}{4.31.6}
\medskip\noindent\textbf{4.31.6}
With the optimised code:
\begin{itemize}
\item One-issue: 4 cycles per iteration (no stalls).
\item Two-issue: 3 cycles per iteration.
\end{itemize}
\[
\text{Speedup} = \frac{4}{3} \approx 1.33
\]
Same as the unoptimised code.  The reordering helped utilisation but did not
reduce the cycle count because the bottleneck is having three ALU/branch
instructions vs.\ one memory instruction per iteration.

\phantomsection\addcontentsline{toc}{subsubsection}{4.31.7}
\medskip\noindent\textbf{4.31.7}
Unroll the loop so each iteration handles two original iterations.  Use
\texttt{\$s5}--\texttt{\$s7} as additional registers:

\begin{verbatim}
    TOP:  lw   $s2, 0($s3)         # iter 1: load
          sub  $s3, $s3, $s4       # iter 1: decrement pointer
          add  $s1, $s2, $s0       # iter 1: accumulate
          lw   $s5, 0($s3)         # iter 2: load (using updated $s3)
          sub  $s3, $s3, $s4       # iter 2: decrement pointer
          add  $s1, $s5, $s0       # iter 2: accumulate
    ENT:  bne  $s0, $s1, TOP       # branch
\end{verbatim}

\noindent Rearranged for the one-issue processor to minimise stalls:

\begin{verbatim}
    TOP:  lw   $s2, 0($s3)         # Memory
          sub  $s3, $s3, $s4       # ALU (fills load-use delay)
          add  $s1, $s2, $s0       # ALU (fwd from lw 2 instr ago)
          lw   $s5, 0($s3)         # Memory (uses updated $s3)
          sub  $s3, $s3, $s4       # ALU (fills load-use delay)
          add  $s1, $s5, $s0       # ALU (fwd from lw 2 instr ago)
    ENT:  bne  $s0, $s1, TOP       # Branch
\end{verbatim}

\noindent One-issue: 7 instructions per unrolled iteration (handling 2
original iterations), with no stalls if forwarding eliminates the load-use
hazard (the \texttt{sub} between \texttt{lw} and \texttt{add} provides the
needed one-cycle gap).  That is 3.5 cycles per original iteration, compared
to 4 cycles without unrolling.

\phantomsection\addcontentsline{toc}{subsubsection}{4.31.8}
\medskip\noindent\textbf{4.31.8}
Two-issue schedule for the unrolled code:

\begin{center}
\begin{tabular}{clll}
\toprule
\textbf{Cycle} & \textbf{Memory slot} & \textbf{ALU/Branch slot} & \textbf{Notes} \\
\midrule
1 & \texttt{lw \$s2, 0(\$s3)} & \texttt{sub \$s3, \$s3, \$s4} & no inter-dep \\
2 & \texttt{nop} & \texttt{add \$s1, \$s2, \$s0} & fwd \texttt{\$s2} \\
3 & \texttt{lw \$s5, 0(\$s3)} & \texttt{sub \$s3, \$s3, \$s4} & no inter-dep \\
4 & \texttt{nop} & \texttt{add \$s1, \$s5, \$s0} & fwd \texttt{\$s5} \\
5 & \texttt{nop} & \texttt{bne \$s0, \$s1, TOP} & fwd \texttt{\$s1} \\
\bottomrule
\end{tabular}
\end{center}

\noindent 5 cycles for 2 original iterations = \textbf{2.5 cycles per
original iteration}.

\phantomsection\addcontentsline{toc}{subsubsection}{4.31.9}
\medskip\noindent\textbf{4.31.9}
\begin{itemize}
\item One-issue (unrolled): 7 cycles for 2 iterations = 3.5 cycles/iteration.
\item Two-issue (unrolled): 5 cycles for 2 iterations = 2.5 cycles/iteration.
\end{itemize}
\[
\text{Speedup} = \frac{3.5}{2.5} = 1.4
\]

\phantomsection\addcontentsline{toc}{subsubsection}{4.31.10}
\medskip\noindent\textbf{4.31.10}
If the two-issue processor can pair two ALU/branch instructions (but not two
memory operations), we can fill the empty ALU slots:

\begin{center}
\begin{tabular}{clll}
\toprule
\textbf{Cycle} & \textbf{Slot 1} & \textbf{Slot 2} & \textbf{Notes} \\
\midrule
1 & \texttt{lw \$s2, 0(\$s3)} & \texttt{sub \$s3, \$s3, \$s4} & mem + ALU \\
2 & \texttt{add \$s1, \$s2, \$s0} & \texttt{lw \$s5, 0(\$s3)} & ALU + mem \\
3 & \texttt{sub \$s3, \$s3, \$s4} & \texttt{add \$s1, \$s5, \$s0} & ALU + ALU, but \texttt{\$s5} from lw in cycle 2 \\
\bottomrule
\end{tabular}
\end{center}

\noindent Wait --- \texttt{add \$s1, \$s5, \$s0} in cycle 3 depends on
\texttt{lw \$s5} from cycle 2 (load-use hazard: need 1 stall).  Rearrange:

\begin{center}
\begin{tabular}{clll}
\toprule
\textbf{Cycle} & \textbf{Slot 1} & \textbf{Slot 2} & \textbf{Notes} \\
\midrule
1 & \texttt{lw \$s2, 0(\$s3)} & \texttt{sub \$s3, \$s3, \$s4} & mem + ALU \\
2 & \texttt{lw \$s5, 0(\$s3)} & \texttt{add \$s1, \$s2, \$s0} & mem + ALU; \texttt{\$s2} fwd OK \\
3 & \texttt{sub \$s3, \$s3, \$s4} & \texttt{add \$s1, \$s5, \$s0} & ALU + ALU; load-use on \texttt{\$s5}! \\
\bottomrule
\end{tabular}
\end{center}

\noindent Still a load-use hazard.  Insert a stall or move \texttt{bne}:

\begin{center}
\begin{tabular}{clll}
\toprule
\textbf{Cycle} & \textbf{Slot 1} & \textbf{Slot 2} & \textbf{Notes} \\
\midrule
1 & \texttt{lw \$s2, 0(\$s3)} & \texttt{sub \$s3, \$s3, \$s4} & mem + ALU \\
2 & \texttt{lw \$s5, 0(\$s3)} & \texttt{add \$s1, \$s2, \$s0} & mem + ALU \\
3 & \texttt{sub \$s3, \$s3, \$s4} & \texttt{nop} & ALU only (wait for \texttt{\$s5}) \\
4 & \texttt{add \$s1, \$s5, \$s0} & \texttt{bne \$s0, \$s1, TOP} & ALU + branch \\
\bottomrule
\end{tabular}
\end{center}

\noindent Wait --- \texttt{add} and \texttt{bne} in cycle 4: \texttt{bne}
reads \texttt{\$s1} which \texttt{add} writes.  They are in the same packet
and one depends on the other --- not allowed.

\begin{center}
\begin{tabular}{clll}
\toprule
\textbf{Cycle} & \textbf{Slot 1} & \textbf{Slot 2} & \textbf{Notes} \\
\midrule
1 & \texttt{lw \$s2, 0(\$s3)} & \texttt{sub \$s3, \$s3, \$s4} & mem + ALU \\
2 & \texttt{lw \$s5, 0(\$s3)} & \texttt{add \$s1, \$s2, \$s0} & mem + ALU \\
3 & \texttt{sub \$s3, \$s3, \$s4} & \texttt{add \$s1, \$s5, \$s0} & ALU + ALU; load-use stall \\
\bottomrule
\end{tabular}
\end{center}

\noindent The load-use hazard on \texttt{\$s5} forces a stall.  Best
achievable schedule:

\begin{center}
\begin{tabular}{clll}
\toprule
\textbf{Cycle} & \textbf{Slot 1} & \textbf{Slot 2} & \textbf{Notes} \\
\midrule
1 & \texttt{lw \$s2, 0(\$s3)} & \texttt{sub \$s3, \$s3, \$s4} & mem + ALU \\
2 & \texttt{lw \$s5, 0(\$s3)} & \texttt{add \$s1, \$s2, \$s0} & mem + ALU \\
3 & \texttt{nop} & \texttt{sub \$s3, \$s3, \$s4} & wait for \texttt{\$s5} \\
4 & \texttt{nop} & \texttt{add \$s1, \$s5, \$s0} & fwd \texttt{\$s5} \\
5 & \texttt{nop} & \texttt{bne \$s0, \$s1, TOP} & fwd \texttt{\$s1} \\
\bottomrule
\end{tabular}
\end{center}

\noindent 5 cycles for 2 iterations = 2.5 cycles per original iteration
(same as before).

Speedup over one-issue:
\[
\text{Speedup} = \frac{3.5}{2.5} = 1.4
\]

\noindent The relaxed pairing constraint did not help here because the
bottleneck is the load-use hazard, not the slot restriction.

%-----------------------------------------------------------------------------
\subsection*{4.32 --- Energy Efficiency Analysis}
\hypertarget{sol-4.32}{}
\markright{4.32}
\addcontentsline{toc}{subsection}{4.32 --- Energy Efficiency Analysis}

Energy consumption per functional unit:

\begin{center}
\begin{tabular}{cccccc}
\toprule
\textbf{I-Mem} & \textbf{Reg Read} & \textbf{Reg Write} & \textbf{ALU} &
\textbf{D-Mem Read} & \textbf{D-Mem Write} \\
\midrule
140\,pJ & 60\,pJ & 80\,pJ & 95\,pJ & 80\,pJ & 55\,pJ \\
\bottomrule
\end{tabular}
\end{center}

\phantomsection\addcontentsline{toc}{subsubsection}{4.32.1}
\medskip\noindent\textbf{4.32.1}
Energy for \texttt{add} in a \textbf{single-cycle} design:

The \texttt{add} instruction activates only the functional units it actually
uses.  Both \texttt{MemRead} and \texttt{MemWrite} are 0 for R-type
instructions, so neither data-memory read nor write energy is consumed:
\[
E_{\text{add, single}} = \underbrace{140}_{\text{I-Mem}}
  + \underbrace{60}_{\text{Reg Read}}
  + \underbrace{95}_{\text{ALU}}
  + \underbrace{80}_{\text{Reg Write}}
  = 375 \text{ pJ}
\]

In a \textbf{pipelined} design, the same functional units are traversed by the
\texttt{add} instruction as it progresses through the five stages.  During the
MEM stage, \texttt{MemRead}~=~0 gates the data memory, so no read energy is
consumed:
\[
E_{\text{add, pipe}} = 140 + 60 + 95 + 80 = 375 \text{ pJ}
\]

\noindent The energy per \texttt{add} instruction is the same in both designs.
(The pipeline saves \emph{time} through parallelism, but each instruction
still visits the same functional units.)

\phantomsection\addcontentsline{toc}{subsubsection}{4.32.2}
\medskip\noindent\textbf{4.32.2}
The worst-case instruction consumes the most energy.  Compare all instruction
types (only counting functional units that are actually activated by the
control signals):

\begin{center}
\begin{tabular}{lccccccl}
\toprule
\textbf{Instr} & \textbf{I-Mem} & \textbf{Reg Rd} & \textbf{ALU} &
\textbf{D-Mem Rd} & \textbf{D-Mem Wr} & \textbf{Reg Wr} & \textbf{Total} \\
\midrule
R-type & 140 & 60 & 95 & --- & --- & 80 & 375 pJ \\
\texttt{lw} & 140 & 60 & 95 & 80 & --- & 80 & 455 pJ \\
\texttt{sw} & 140 & 60 & 95 & --- & 55 & --- & 350 pJ \\
\texttt{beq} & 140 & 60 & 95 & --- & --- & --- & 295 pJ \\
\bottomrule
\end{tabular}
\end{center}

\noindent The worst-case instruction is \texttt{lw} at $\boxed{455 \text{ pJ}}$.
It is the only instruction that activates every major functional unit: I-Mem,
Reg Read, ALU (address calculation), D-Mem Read, and Reg Write.  R-type
instructions are second at 375\,pJ (no data memory access).

\phantomsection\addcontentsline{toc}{subsubsection}{4.32.3}
\medskip\noindent\textbf{4.32.3}
The \texttt{MemRead} signal already gates the data memory read port, so
non-load instructions do not waste D-Mem energy in the current design.
The main opportunity to reduce \texttt{lw} energy is to use the
\textbf{no-offset load/store} design (from Exercise~4.15), which removes the
ALU from the load path (the register value goes directly to the memory address
input):
\[
E_{\text{lw, new}} = 140 + 60 + 80 + 80 = 360 \text{ pJ}
\]
\[
\text{Reduction} = \frac{455 - 360}{455} = \frac{95}{455} \approx 20.9\%
\]

\noindent The 95\,pJ savings comes entirely from eliminating the ALU from the
load critical path.

\phantomsection\addcontentsline{toc}{subsubsection}{4.32.4}
\medskip\noindent\textbf{4.32.4}
Other instructions that benefit from the no-offset load/store change:

\begin{itemize}
\item \textbf{\texttt{sw}} also benefits because it no longer needs the ALU to
      compute the address:
      $E_{\text{sw, new}} = 140 + 60 + 55 = 255$\,pJ (down from 350\,pJ,
      a 27.1\% reduction).

\item \textbf{R-type, \texttt{beq}, and other non-memory instructions} are
      unaffected by this change --- they still use the ALU for their primary
      computation.
\end{itemize}

\phantomsection\addcontentsline{toc}{subsubsection}{4.32.5}
\medskip\noindent\textbf{4.32.5}
The no-offset load/store change (from Exercise 4.15) has the following
performance effects on a pipelined CPU:

\begin{itemize}
\item \textbf{Positive:} The clock cycle time may decrease because load/store
      no longer uses both the ALU and data memory in sequence.  The MEM and
      EX stages could potentially overlap (4-stage pipeline).

\item \textbf{Negative:} The instruction count \emph{increases} because every
      load/store with a non-zero offset now requires an additional \texttt{add}
      instruction to pre-compute the address.  If 35\% of instructions are
      loads/stores (from the Exercise 4.3 mix: 25\% + 10\%), many will need
      extra instructions.

\item \textbf{Net effect:}  As shown in Exercise 4.15, the increase in
      instruction count typically outweighs the decrease in clock cycle time,
      resulting in a net \emph{slowdown} in most programs.  The energy savings
      must be weighed against this performance degradation.
\end{itemize}

\phantomsection\addcontentsline{toc}{subsubsection}{4.32.6}
\medskip\noindent\textbf{4.32.6}
\textbf{Can we eliminate the MemRead control signal?}

Yes.  If data memory is readable every cycle, the processor still functions
correctly because:
\begin{itemize}
\item For \texttt{lw}: the data memory output is selected by the MemtoReg mux
      (MemtoReg = 1) and written to the register file.  Behaviour unchanged.
\item For all other instructions: the data memory still produces an output,
      but the MemtoReg mux selects the ALU result (MemtoReg = 0), so the
      memory output is simply \emph{ignored}.  No incorrect data is written.
\end{itemize}

\noindent Effects:

\textbf{Clock frequency:}  No change.  Data memory was already on the critical
path for \texttt{lw} instructions, and its latency does not increase.  For
non-load instructions, the memory read happens in parallel with the ALU (or
is an unused output), so it does not lengthen the critical path.

\textbf{Energy consumption:}  With 25\% loads, the remaining 75\% of
instructions now waste energy on unnecessary D-Mem reads:
\[
\text{Extra energy per instruction (avg)} = 0.75 \times 80 = 60 \text{ pJ}
\]

\noindent For a program of $N$ instructions, the total wasted energy is
$60N$\,pJ.  If the average instruction energy is approximately 380\,pJ
(weighted average with MemRead gating), this represents about a 16\% increase
in energy consumption.  The trade-off: simpler control logic (one fewer signal)
at the cost of higher energy.

%-----------------------------------------------------------------------------
\subsection*{4.33 --- Stuck-at Fault Testing}
\hypertarget{sol-4.33}{}
\markright{4.33}
\addcontentsline{toc}{subsection}{4.33 --- Stuck-at Fault Testing}

All problems refer to \textbf{bit 0 of the Write Register input} of the
Register File, unless otherwise stated.

\phantomsection\addcontentsline{toc}{subsubsection}{4.33.1}
\medskip\noindent\textbf{4.33.1}
\textbf{Test for stuck-at-0 on bit 0 of Write Register.}

If bit 0 is stuck at 0, any write to an odd-numbered register (bit 0 = 1)
will instead be directed to the corresponding even-numbered register
(bit 0 forced to 0).  For example, a write to register 1 (\texttt{00001})
goes to register 0 (\texttt{00000}).

\textbf{Test procedure:}
\begin{enumerate}
\item Initialise register \texttt{\$1} to a known value (e.g., 0).
\item Execute: \texttt{addi \$1, \$zero, 42}
      \quad (should write 42 to register 1).
\item Read register \texttt{\$1}.
\end{enumerate}

\textbf{Expected results:}
\begin{itemize}
\item \textbf{No fault:} \texttt{\$1} = 42. \checkmark
\item \textbf{Stuck-at-0:} The write targets register 0 (\texttt{00000})
      instead of register 1 (\texttt{00001}).  Since \texttt{\$0} is
      hardwired to zero, the write has no effect.  \texttt{\$1} retains its
      initial value of 0.  Detecting \texttt{\$1} $\neq$ 42 reveals the fault.
\end{itemize}

\noindent More generally, any R-type or I-type instruction that writes to
an odd-numbered register can serve as a test.  After execution, verify the
odd register holds the expected result.

\phantomsection\addcontentsline{toc}{subsubsection}{4.33.2}
\medskip\noindent\textbf{4.33.2}
\textbf{Test for stuck-at-1 on bit 0 of Write Register.}

If bit 0 is stuck at 1, any write to an even-numbered register (bit 0 = 0)
will instead go to the corresponding odd-numbered register (bit 0 forced
to 1).  For example, a write to register 2 (\texttt{00010}) goes to
register 3 (\texttt{00011}).

\textbf{Test procedure:}
\begin{enumerate}
\item Initialise register \texttt{\$2} to 0 and register \texttt{\$3} to 99.
\item Execute: \texttt{addi \$2, \$zero, 42}
      \quad (should write 42 to register 2).
\item Read registers \texttt{\$2} and \texttt{\$3}.
\end{enumerate}

\textbf{Expected results:}
\begin{itemize}
\item \textbf{No fault:} \texttt{\$2} = 42, \texttt{\$3} = 99. \checkmark
\item \textbf{Stuck-at-1:} The write goes to register 3 instead of register 2.
      \texttt{\$2} remains 0, \texttt{\$3} becomes 42 (instead of 99).
      Detecting \texttt{\$2} $\neq$ 42 or \texttt{\$3} $\neq$ 99 reveals
      the fault.
\end{itemize}

\phantomsection\addcontentsline{toc}{subsubsection}{4.33.3}
\medskip\noindent\textbf{4.33.3}
\textbf{Test for RegWrite stuck-at-0.}

If RegWrite is stuck at 0, \emph{no instruction} can write to the register
file---every register write is suppressed.

\textbf{Test procedure:}
\begin{enumerate}
\item Initialise register \texttt{\$1} to 0.
\item Execute: \texttt{addi \$1, \$zero, 7}
      \quad (should write 7 to \texttt{\$1}).
\item Read register \texttt{\$1}.
\end{enumerate}

\textbf{Expected results:}
\begin{itemize}
\item \textbf{No fault:} \texttt{\$1} = 7. \checkmark
\item \textbf{Stuck-at-0:} RegWrite is never asserted, so the register file
      ignores the write.  \texttt{\$1} remains 0.  Detecting
      \texttt{\$1} $\neq$ 7 reveals the fault.
\end{itemize}

\noindent Any instruction that writes a register (\texttt{add}, \texttt{sub},
\texttt{lw}, \texttt{addi}, etc.) can serve as a test.

\phantomsection\addcontentsline{toc}{subsubsection}{4.33.4}
\medskip\noindent\textbf{4.33.4}
\textbf{Test for RegWrite stuck-at-1.}

If RegWrite is stuck at 1, \emph{every instruction} writes to the register
file---including instructions that should not (e.g., \texttt{sw},
\texttt{beq}).  The value written and the destination register depend on the
datapath signals:
\begin{itemize}
\item The write-register input comes from the RegDst mux (either \texttt{rd}
      or \texttt{rt} field).
\item The write-data input comes from the MemtoReg mux (either ALU result or
      memory data).
\end{itemize}

\textbf{Test procedure:}
\begin{enumerate}
\item Initialise all registers to known values (e.g., \texttt{\$1} = 100,
      \texttt{\$2} = 200, etc.).
\item Execute an instruction that should \emph{not} write to the register
      file.  For example:

\begin{verbatim}
    sw $1, 0($2)     # Store; RegWrite should be 0
\end{verbatim}

\item Read all registers and compare to their initial values.
\end{enumerate}

\textbf{Expected results:}
\begin{itemize}
\item \textbf{No fault:} All registers retain their initial values. \checkmark
\item \textbf{Stuck-at-1:} The \texttt{sw} instruction writes a value to
      the register file.  For \texttt{sw}, the RegDst mux output selects a
      destination register (the \texttt{rd} field, which for \texttt{sw} is
      bits [15:11] of the instruction encoding, typically a garbage destination).
      The MemtoReg mux selects the ALU result (the computed memory address
      $\texttt{\$2} + 0$).  Some register is overwritten with the address
      value.  Detecting any register change reveals the fault.
\end{itemize}

\noindent A \texttt{beq} instruction is equally effective: it should not write
any register, so any register modification after executing \texttt{beq}
indicates a stuck-at-1 fault on RegWrite.
